\documentclass[11pt]{article}

    \usepackage[breakable]{tcolorbox}
    \usepackage{parskip} % Stop auto-indenting (to mimic markdown behaviour)
    

    % Basic figure setup, for now with no caption control since it's done
    % automatically by Pandoc (which extracts ![](path) syntax from Markdown).
    \usepackage{graphicx}
    % Keep aspect ratio if custom image width or height is specified
    \setkeys{Gin}{keepaspectratio}
    % Maintain compatibility with old templates. Remove in nbconvert 6.0
    \let\Oldincludegraphics\includegraphics
    % Ensure that by default, figures have no caption (until we provide a
    % proper Figure object with a Caption API and a way to capture that
    % in the conversion process - todo).
    \usepackage{caption}
    \DeclareCaptionFormat{nocaption}{}
    \captionsetup{format=nocaption,aboveskip=0pt,belowskip=0pt}

    \usepackage{float}
    \floatplacement{figure}{H} % forces figures to be placed at the correct location
    \usepackage{xcolor} % Allow colors to be defined
    \usepackage{enumerate} % Needed for markdown enumerations to work
    \usepackage{geometry} % Used to adjust the document margins
    \usepackage{amsmath} % Equations
    \usepackage{amssymb} % Equations
    \usepackage{textcomp} % defines textquotesingle
    % Hack from http://tex.stackexchange.com/a/47451/13684:
    \AtBeginDocument{%
        \def\PYZsq{\textquotesingle}% Upright quotes in Pygmentized code
    }
    \usepackage{upquote} % Upright quotes for verbatim code
    \usepackage{eurosym} % defines \euro

    \usepackage{iftex}
    \ifPDFTeX
        \usepackage[T1]{fontenc}
        \IfFileExists{alphabeta.sty}{
              \usepackage{alphabeta}
          }{
              \usepackage[mathletters]{ucs}
              \usepackage[utf8x]{inputenc}
          }
    \else
        \usepackage{fontspec}
        \usepackage{unicode-math}
    \fi

    \usepackage{fancyvrb} % verbatim replacement that allows latex
    \usepackage{grffile} % extends the file name processing of package graphics
                         % to support a larger range
    \makeatletter % fix for old versions of grffile with XeLaTeX
    \@ifpackagelater{grffile}{2019/11/01}
    {
      % Do nothing on new versions
    }
    {
      \def\Gread@@xetex#1{%
        \IfFileExists{"\Gin@base".bb}%
        {\Gread@eps{\Gin@base.bb}}%
        {\Gread@@xetex@aux#1}%
      }
    }
    \makeatother
    \usepackage[Export]{adjustbox} % Used to constrain images to a maximum size
    \adjustboxset{max size={0.9\linewidth}{0.9\paperheight}}

    % The hyperref package gives us a pdf with properly built
    % internal navigation ('pdf bookmarks' for the table of contents,
    % internal cross-reference links, web links for URLs, etc.)
    \usepackage{hyperref}
    % The default LaTeX title has an obnoxious amount of whitespace. By default,
    % titling removes some of it. It also provides customization options.
    \usepackage{titling}
    \usepackage{longtable} % longtable support required by pandoc >1.10
    \usepackage{booktabs}  % table support for pandoc > 1.12.2
    \usepackage{array}     % table support for pandoc >= 2.11.3
    \usepackage{calc}      % table minipage width calculation for pandoc >= 2.11.1
    \usepackage[inline]{enumitem} % IRkernel/repr support (it uses the enumerate* environment)
    \usepackage[normalem]{ulem} % ulem is needed to support strikethroughs (\sout)
                                % normalem makes italics be italics, not underlines
    \usepackage{soul}      % strikethrough (\st) support for pandoc >= 3.0.0
    \usepackage{mathrsfs}
    

    
    % Colors for the hyperref package
    \definecolor{urlcolor}{rgb}{0,.145,.698}
    \definecolor{linkcolor}{rgb}{.71,0.21,0.01}
    \definecolor{citecolor}{rgb}{.12,.54,.11}

    % ANSI colors
    \definecolor{ansi-black}{HTML}{3E424D}
    \definecolor{ansi-black-intense}{HTML}{282C36}
    \definecolor{ansi-red}{HTML}{E75C58}
    \definecolor{ansi-red-intense}{HTML}{B22B31}
    \definecolor{ansi-green}{HTML}{00A250}
    \definecolor{ansi-green-intense}{HTML}{007427}
    \definecolor{ansi-yellow}{HTML}{DDB62B}
    \definecolor{ansi-yellow-intense}{HTML}{B27D12}
    \definecolor{ansi-blue}{HTML}{208FFB}
    \definecolor{ansi-blue-intense}{HTML}{0065CA}
    \definecolor{ansi-magenta}{HTML}{D160C4}
    \definecolor{ansi-magenta-intense}{HTML}{A03196}
    \definecolor{ansi-cyan}{HTML}{60C6C8}
    \definecolor{ansi-cyan-intense}{HTML}{258F8F}
    \definecolor{ansi-white}{HTML}{C5C1B4}
    \definecolor{ansi-white-intense}{HTML}{A1A6B2}
    \definecolor{ansi-default-inverse-fg}{HTML}{FFFFFF}
    \definecolor{ansi-default-inverse-bg}{HTML}{000000}

    % common color for the border for error outputs.
    \definecolor{outerrorbackground}{HTML}{FFDFDF}

    % commands and environments needed by pandoc snippets
    % extracted from the output of `pandoc -s`
    \providecommand{\tightlist}{%
      \setlength{\itemsep}{0pt}\setlength{\parskip}{0pt}}
    \DefineVerbatimEnvironment{Highlighting}{Verbatim}{commandchars=\\\{\}}
    % Add ',fontsize=\small' for more characters per line
    \newenvironment{Shaded}{}{}
    \newcommand{\KeywordTok}[1]{\textcolor[rgb]{0.00,0.44,0.13}{\textbf{{#1}}}}
    \newcommand{\DataTypeTok}[1]{\textcolor[rgb]{0.56,0.13,0.00}{{#1}}}
    \newcommand{\DecValTok}[1]{\textcolor[rgb]{0.25,0.63,0.44}{{#1}}}
    \newcommand{\BaseNTok}[1]{\textcolor[rgb]{0.25,0.63,0.44}{{#1}}}
    \newcommand{\FloatTok}[1]{\textcolor[rgb]{0.25,0.63,0.44}{{#1}}}
    \newcommand{\CharTok}[1]{\textcolor[rgb]{0.25,0.44,0.63}{{#1}}}
    \newcommand{\StringTok}[1]{\textcolor[rgb]{0.25,0.44,0.63}{{#1}}}
    \newcommand{\CommentTok}[1]{\textcolor[rgb]{0.38,0.63,0.69}{\textit{{#1}}}}
    \newcommand{\OtherTok}[1]{\textcolor[rgb]{0.00,0.44,0.13}{{#1}}}
    \newcommand{\AlertTok}[1]{\textcolor[rgb]{1.00,0.00,0.00}{\textbf{{#1}}}}
    \newcommand{\FunctionTok}[1]{\textcolor[rgb]{0.02,0.16,0.49}{{#1}}}
    \newcommand{\RegionMarkerTok}[1]{{#1}}
    \newcommand{\ErrorTok}[1]{\textcolor[rgb]{1.00,0.00,0.00}{\textbf{{#1}}}}
    \newcommand{\NormalTok}[1]{{#1}}

    % Additional commands for more recent versions of Pandoc
    \newcommand{\ConstantTok}[1]{\textcolor[rgb]{0.53,0.00,0.00}{{#1}}}
    \newcommand{\SpecialCharTok}[1]{\textcolor[rgb]{0.25,0.44,0.63}{{#1}}}
    \newcommand{\VerbatimStringTok}[1]{\textcolor[rgb]{0.25,0.44,0.63}{{#1}}}
    \newcommand{\SpecialStringTok}[1]{\textcolor[rgb]{0.73,0.40,0.53}{{#1}}}
    \newcommand{\ImportTok}[1]{{#1}}
    \newcommand{\DocumentationTok}[1]{\textcolor[rgb]{0.73,0.13,0.13}{\textit{{#1}}}}
    \newcommand{\AnnotationTok}[1]{\textcolor[rgb]{0.38,0.63,0.69}{\textbf{\textit{{#1}}}}}
    \newcommand{\CommentVarTok}[1]{\textcolor[rgb]{0.38,0.63,0.69}{\textbf{\textit{{#1}}}}}
    \newcommand{\VariableTok}[1]{\textcolor[rgb]{0.10,0.09,0.49}{{#1}}}
    \newcommand{\ControlFlowTok}[1]{\textcolor[rgb]{0.00,0.44,0.13}{\textbf{{#1}}}}
    \newcommand{\OperatorTok}[1]{\textcolor[rgb]{0.40,0.40,0.40}{{#1}}}
    \newcommand{\BuiltInTok}[1]{{#1}}
    \newcommand{\ExtensionTok}[1]{{#1}}
    \newcommand{\PreprocessorTok}[1]{\textcolor[rgb]{0.74,0.48,0.00}{{#1}}}
    \newcommand{\AttributeTok}[1]{\textcolor[rgb]{0.49,0.56,0.16}{{#1}}}
    \newcommand{\InformationTok}[1]{\textcolor[rgb]{0.38,0.63,0.69}{\textbf{\textit{{#1}}}}}
    \newcommand{\WarningTok}[1]{\textcolor[rgb]{0.38,0.63,0.69}{\textbf{\textit{{#1}}}}}


    % Define a nice break command that doesn't care if a line doesn't already
    % exist.
    \def\br{\hspace*{\fill} \\* }
    % Math Jax compatibility definitions
    \def\gt{>}
    \def\lt{<}
    \let\Oldtex\TeX
    \let\Oldlatex\LaTeX
    \renewcommand{\TeX}{\textrm{\Oldtex}}
    \renewcommand{\LaTeX}{\textrm{\Oldlatex}}
    % Document parameters
    % Document title
    \title{ntt}
    
    
    
    
    
    
    
% Pygments definitions
\makeatletter
\def\PY@reset{\let\PY@it=\relax \let\PY@bf=\relax%
    \let\PY@ul=\relax \let\PY@tc=\relax%
    \let\PY@bc=\relax \let\PY@ff=\relax}
\def\PY@tok#1{\csname PY@tok@#1\endcsname}
\def\PY@toks#1+{\ifx\relax#1\empty\else%
    \PY@tok{#1}\expandafter\PY@toks\fi}
\def\PY@do#1{\PY@bc{\PY@tc{\PY@ul{%
    \PY@it{\PY@bf{\PY@ff{#1}}}}}}}
\def\PY#1#2{\PY@reset\PY@toks#1+\relax+\PY@do{#2}}

\@namedef{PY@tok@w}{\def\PY@tc##1{\textcolor[rgb]{0.73,0.73,0.73}{##1}}}
\@namedef{PY@tok@c}{\let\PY@it=\textit\def\PY@tc##1{\textcolor[rgb]{0.24,0.48,0.48}{##1}}}
\@namedef{PY@tok@cp}{\def\PY@tc##1{\textcolor[rgb]{0.61,0.40,0.00}{##1}}}
\@namedef{PY@tok@k}{\let\PY@bf=\textbf\def\PY@tc##1{\textcolor[rgb]{0.00,0.50,0.00}{##1}}}
\@namedef{PY@tok@kp}{\def\PY@tc##1{\textcolor[rgb]{0.00,0.50,0.00}{##1}}}
\@namedef{PY@tok@kt}{\def\PY@tc##1{\textcolor[rgb]{0.69,0.00,0.25}{##1}}}
\@namedef{PY@tok@o}{\def\PY@tc##1{\textcolor[rgb]{0.40,0.40,0.40}{##1}}}
\@namedef{PY@tok@ow}{\let\PY@bf=\textbf\def\PY@tc##1{\textcolor[rgb]{0.67,0.13,1.00}{##1}}}
\@namedef{PY@tok@nb}{\def\PY@tc##1{\textcolor[rgb]{0.00,0.50,0.00}{##1}}}
\@namedef{PY@tok@nf}{\def\PY@tc##1{\textcolor[rgb]{0.00,0.00,1.00}{##1}}}
\@namedef{PY@tok@nc}{\let\PY@bf=\textbf\def\PY@tc##1{\textcolor[rgb]{0.00,0.00,1.00}{##1}}}
\@namedef{PY@tok@nn}{\let\PY@bf=\textbf\def\PY@tc##1{\textcolor[rgb]{0.00,0.00,1.00}{##1}}}
\@namedef{PY@tok@ne}{\let\PY@bf=\textbf\def\PY@tc##1{\textcolor[rgb]{0.80,0.25,0.22}{##1}}}
\@namedef{PY@tok@nv}{\def\PY@tc##1{\textcolor[rgb]{0.10,0.09,0.49}{##1}}}
\@namedef{PY@tok@no}{\def\PY@tc##1{\textcolor[rgb]{0.53,0.00,0.00}{##1}}}
\@namedef{PY@tok@nl}{\def\PY@tc##1{\textcolor[rgb]{0.46,0.46,0.00}{##1}}}
\@namedef{PY@tok@ni}{\let\PY@bf=\textbf\def\PY@tc##1{\textcolor[rgb]{0.44,0.44,0.44}{##1}}}
\@namedef{PY@tok@na}{\def\PY@tc##1{\textcolor[rgb]{0.41,0.47,0.13}{##1}}}
\@namedef{PY@tok@nt}{\let\PY@bf=\textbf\def\PY@tc##1{\textcolor[rgb]{0.00,0.50,0.00}{##1}}}
\@namedef{PY@tok@nd}{\def\PY@tc##1{\textcolor[rgb]{0.67,0.13,1.00}{##1}}}
\@namedef{PY@tok@s}{\def\PY@tc##1{\textcolor[rgb]{0.73,0.13,0.13}{##1}}}
\@namedef{PY@tok@sd}{\let\PY@it=\textit\def\PY@tc##1{\textcolor[rgb]{0.73,0.13,0.13}{##1}}}
\@namedef{PY@tok@si}{\let\PY@bf=\textbf\def\PY@tc##1{\textcolor[rgb]{0.64,0.35,0.47}{##1}}}
\@namedef{PY@tok@se}{\let\PY@bf=\textbf\def\PY@tc##1{\textcolor[rgb]{0.67,0.36,0.12}{##1}}}
\@namedef{PY@tok@sr}{\def\PY@tc##1{\textcolor[rgb]{0.64,0.35,0.47}{##1}}}
\@namedef{PY@tok@ss}{\def\PY@tc##1{\textcolor[rgb]{0.10,0.09,0.49}{##1}}}
\@namedef{PY@tok@sx}{\def\PY@tc##1{\textcolor[rgb]{0.00,0.50,0.00}{##1}}}
\@namedef{PY@tok@m}{\def\PY@tc##1{\textcolor[rgb]{0.40,0.40,0.40}{##1}}}
\@namedef{PY@tok@gh}{\let\PY@bf=\textbf\def\PY@tc##1{\textcolor[rgb]{0.00,0.00,0.50}{##1}}}
\@namedef{PY@tok@gu}{\let\PY@bf=\textbf\def\PY@tc##1{\textcolor[rgb]{0.50,0.00,0.50}{##1}}}
\@namedef{PY@tok@gd}{\def\PY@tc##1{\textcolor[rgb]{0.63,0.00,0.00}{##1}}}
\@namedef{PY@tok@gi}{\def\PY@tc##1{\textcolor[rgb]{0.00,0.52,0.00}{##1}}}
\@namedef{PY@tok@gr}{\def\PY@tc##1{\textcolor[rgb]{0.89,0.00,0.00}{##1}}}
\@namedef{PY@tok@ge}{\let\PY@it=\textit}
\@namedef{PY@tok@gs}{\let\PY@bf=\textbf}
\@namedef{PY@tok@gp}{\let\PY@bf=\textbf\def\PY@tc##1{\textcolor[rgb]{0.00,0.00,0.50}{##1}}}
\@namedef{PY@tok@go}{\def\PY@tc##1{\textcolor[rgb]{0.44,0.44,0.44}{##1}}}
\@namedef{PY@tok@gt}{\def\PY@tc##1{\textcolor[rgb]{0.00,0.27,0.87}{##1}}}
\@namedef{PY@tok@err}{\def\PY@bc##1{{\setlength{\fboxsep}{\string -\fboxrule}\fcolorbox[rgb]{1.00,0.00,0.00}{1,1,1}{\strut ##1}}}}
\@namedef{PY@tok@kc}{\let\PY@bf=\textbf\def\PY@tc##1{\textcolor[rgb]{0.00,0.50,0.00}{##1}}}
\@namedef{PY@tok@kd}{\let\PY@bf=\textbf\def\PY@tc##1{\textcolor[rgb]{0.00,0.50,0.00}{##1}}}
\@namedef{PY@tok@kn}{\let\PY@bf=\textbf\def\PY@tc##1{\textcolor[rgb]{0.00,0.50,0.00}{##1}}}
\@namedef{PY@tok@kr}{\let\PY@bf=\textbf\def\PY@tc##1{\textcolor[rgb]{0.00,0.50,0.00}{##1}}}
\@namedef{PY@tok@bp}{\def\PY@tc##1{\textcolor[rgb]{0.00,0.50,0.00}{##1}}}
\@namedef{PY@tok@fm}{\def\PY@tc##1{\textcolor[rgb]{0.00,0.00,1.00}{##1}}}
\@namedef{PY@tok@vc}{\def\PY@tc##1{\textcolor[rgb]{0.10,0.09,0.49}{##1}}}
\@namedef{PY@tok@vg}{\def\PY@tc##1{\textcolor[rgb]{0.10,0.09,0.49}{##1}}}
\@namedef{PY@tok@vi}{\def\PY@tc##1{\textcolor[rgb]{0.10,0.09,0.49}{##1}}}
\@namedef{PY@tok@vm}{\def\PY@tc##1{\textcolor[rgb]{0.10,0.09,0.49}{##1}}}
\@namedef{PY@tok@sa}{\def\PY@tc##1{\textcolor[rgb]{0.73,0.13,0.13}{##1}}}
\@namedef{PY@tok@sb}{\def\PY@tc##1{\textcolor[rgb]{0.73,0.13,0.13}{##1}}}
\@namedef{PY@tok@sc}{\def\PY@tc##1{\textcolor[rgb]{0.73,0.13,0.13}{##1}}}
\@namedef{PY@tok@dl}{\def\PY@tc##1{\textcolor[rgb]{0.73,0.13,0.13}{##1}}}
\@namedef{PY@tok@s2}{\def\PY@tc##1{\textcolor[rgb]{0.73,0.13,0.13}{##1}}}
\@namedef{PY@tok@sh}{\def\PY@tc##1{\textcolor[rgb]{0.73,0.13,0.13}{##1}}}
\@namedef{PY@tok@s1}{\def\PY@tc##1{\textcolor[rgb]{0.73,0.13,0.13}{##1}}}
\@namedef{PY@tok@mb}{\def\PY@tc##1{\textcolor[rgb]{0.40,0.40,0.40}{##1}}}
\@namedef{PY@tok@mf}{\def\PY@tc##1{\textcolor[rgb]{0.40,0.40,0.40}{##1}}}
\@namedef{PY@tok@mh}{\def\PY@tc##1{\textcolor[rgb]{0.40,0.40,0.40}{##1}}}
\@namedef{PY@tok@mi}{\def\PY@tc##1{\textcolor[rgb]{0.40,0.40,0.40}{##1}}}
\@namedef{PY@tok@il}{\def\PY@tc##1{\textcolor[rgb]{0.40,0.40,0.40}{##1}}}
\@namedef{PY@tok@mo}{\def\PY@tc##1{\textcolor[rgb]{0.40,0.40,0.40}{##1}}}
\@namedef{PY@tok@ch}{\let\PY@it=\textit\def\PY@tc##1{\textcolor[rgb]{0.24,0.48,0.48}{##1}}}
\@namedef{PY@tok@cm}{\let\PY@it=\textit\def\PY@tc##1{\textcolor[rgb]{0.24,0.48,0.48}{##1}}}
\@namedef{PY@tok@cpf}{\let\PY@it=\textit\def\PY@tc##1{\textcolor[rgb]{0.24,0.48,0.48}{##1}}}
\@namedef{PY@tok@c1}{\let\PY@it=\textit\def\PY@tc##1{\textcolor[rgb]{0.24,0.48,0.48}{##1}}}
\@namedef{PY@tok@cs}{\let\PY@it=\textit\def\PY@tc##1{\textcolor[rgb]{0.24,0.48,0.48}{##1}}}

\def\PYZbs{\char`\\}
\def\PYZus{\char`\_}
\def\PYZob{\char`\{}
\def\PYZcb{\char`\}}
\def\PYZca{\char`\^}
\def\PYZam{\char`\&}
\def\PYZlt{\char`\<}
\def\PYZgt{\char`\>}
\def\PYZsh{\char`\#}
\def\PYZpc{\char`\%}
\def\PYZdl{\char`\$}
\def\PYZhy{\char`\-}
\def\PYZsq{\char`\'}
\def\PYZdq{\char`\"}
\def\PYZti{\char`\~}
% for compatibility with earlier versions
\def\PYZat{@}
\def\PYZlb{[}
\def\PYZrb{]}
\makeatother


    % For linebreaks inside Verbatim environment from package fancyvrb.
    \makeatletter
        \newbox\Wrappedcontinuationbox
        \newbox\Wrappedvisiblespacebox
        \newcommand*\Wrappedvisiblespace {\textcolor{red}{\textvisiblespace}}
        \newcommand*\Wrappedcontinuationsymbol {\textcolor{red}{\llap{\tiny$\m@th\hookrightarrow$}}}
        \newcommand*\Wrappedcontinuationindent {3ex }
        \newcommand*\Wrappedafterbreak {\kern\Wrappedcontinuationindent\copy\Wrappedcontinuationbox}
        % Take advantage of the already applied Pygments mark-up to insert
        % potential linebreaks for TeX processing.
        %        {, <, #, %, $, ' and ": go to next line.
        %        _, }, ^, &, >, - and ~: stay at end of broken line.
        % Use of \textquotesingle for straight quote.
        \newcommand*\Wrappedbreaksatspecials {%
            \def\PYGZus{\discretionary{\char`\_}{\Wrappedafterbreak}{\char`\_}}%
            \def\PYGZob{\discretionary{}{\Wrappedafterbreak\char`\{}{\char`\{}}%
            \def\PYGZcb{\discretionary{\char`\}}{\Wrappedafterbreak}{\char`\}}}%
            \def\PYGZca{\discretionary{\char`\^}{\Wrappedafterbreak}{\char`\^}}%
            \def\PYGZam{\discretionary{\char`\&}{\Wrappedafterbreak}{\char`\&}}%
            \def\PYGZlt{\discretionary{}{\Wrappedafterbreak\char`\<}{\char`\<}}%
            \def\PYGZgt{\discretionary{\char`\>}{\Wrappedafterbreak}{\char`\>}}%
            \def\PYGZsh{\discretionary{}{\Wrappedafterbreak\char`\#}{\char`\#}}%
            \def\PYGZpc{\discretionary{}{\Wrappedafterbreak\char`\%}{\char`\%}}%
            \def\PYGZdl{\discretionary{}{\Wrappedafterbreak\char`\$}{\char`\$}}%
            \def\PYGZhy{\discretionary{\char`\-}{\Wrappedafterbreak}{\char`\-}}%
            \def\PYGZsq{\discretionary{}{\Wrappedafterbreak\textquotesingle}{\textquotesingle}}%
            \def\PYGZdq{\discretionary{}{\Wrappedafterbreak\char`\"}{\char`\"}}%
            \def\PYGZti{\discretionary{\char`\~}{\Wrappedafterbreak}{\char`\~}}%
        }
        % Some characters . , ; ? ! / are not pygmentized.
        % This macro makes them "active" and they will insert potential linebreaks
        \newcommand*\Wrappedbreaksatpunct {%
            \lccode`\~`\.\lowercase{\def~}{\discretionary{\hbox{\char`\.}}{\Wrappedafterbreak}{\hbox{\char`\.}}}%
            \lccode`\~`\,\lowercase{\def~}{\discretionary{\hbox{\char`\,}}{\Wrappedafterbreak}{\hbox{\char`\,}}}%
            \lccode`\~`\;\lowercase{\def~}{\discretionary{\hbox{\char`\;}}{\Wrappedafterbreak}{\hbox{\char`\;}}}%
            \lccode`\~`\:\lowercase{\def~}{\discretionary{\hbox{\char`\:}}{\Wrappedafterbreak}{\hbox{\char`\:}}}%
            \lccode`\~`\?\lowercase{\def~}{\discretionary{\hbox{\char`\?}}{\Wrappedafterbreak}{\hbox{\char`\?}}}%
            \lccode`\~`\!\lowercase{\def~}{\discretionary{\hbox{\char`\!}}{\Wrappedafterbreak}{\hbox{\char`\!}}}%
            \lccode`\~`\/\lowercase{\def~}{\discretionary{\hbox{\char`\/}}{\Wrappedafterbreak}{\hbox{\char`\/}}}%
            \catcode`\.\active
            \catcode`\,\active
            \catcode`\;\active
            \catcode`\:\active
            \catcode`\?\active
            \catcode`\!\active
            \catcode`\/\active
            \lccode`\~`\~
        }
    \makeatother

    \let\OriginalVerbatim=\Verbatim
    \makeatletter
    \renewcommand{\Verbatim}[1][1]{%
        %\parskip\z@skip
        \sbox\Wrappedcontinuationbox {\Wrappedcontinuationsymbol}%
        \sbox\Wrappedvisiblespacebox {\FV@SetupFont\Wrappedvisiblespace}%
        \def\FancyVerbFormatLine ##1{\hsize\linewidth
            \vtop{\raggedright\hyphenpenalty\z@\exhyphenpenalty\z@
                \doublehyphendemerits\z@\finalhyphendemerits\z@
                \strut ##1\strut}%
        }%
        % If the linebreak is at a space, the latter will be displayed as visible
        % space at end of first line, and a continuation symbol starts next line.
        % Stretch/shrink are however usually zero for typewriter font.
        \def\FV@Space {%
            \nobreak\hskip\z@ plus\fontdimen3\font minus\fontdimen4\font
            \discretionary{\copy\Wrappedvisiblespacebox}{\Wrappedafterbreak}
            {\kern\fontdimen2\font}%
        }%

        % Allow breaks at special characters using \PYG... macros.
        \Wrappedbreaksatspecials
        % Breaks at punctuation characters . , ; ? ! and / need catcode=\active
        \OriginalVerbatim[#1,codes*=\Wrappedbreaksatpunct]%
    }
    \makeatother

    % Exact colors from NB
    \definecolor{incolor}{HTML}{303F9F}
    \definecolor{outcolor}{HTML}{D84315}
    \definecolor{cellborder}{HTML}{CFCFCF}
    \definecolor{cellbackground}{HTML}{F7F7F7}

    % prompt
    \makeatletter
    \newcommand{\boxspacing}{\kern\kvtcb@left@rule\kern\kvtcb@boxsep}
    \makeatother
    \newcommand{\prompt}[4]{
        {\ttfamily\llap{{\color{#2}[#3]:\hspace{3pt}#4}}\vspace{-\baselineskip}}
    }
    

    
    % Prevent overflowing lines due to hard-to-break entities
    \sloppy
    % Setup hyperref package
    \hypersetup{
      breaklinks=true,  % so long urls are correctly broken across lines
      colorlinks=true,
      urlcolor=urlcolor,
      linkcolor=linkcolor,
      citecolor=citecolor,
      }
    % Slightly bigger margins than the latex defaults
    
    \geometry{verbose,tmargin=1in,bmargin=1in,lmargin=1in,rmargin=1in}
    
    

\begin{document}
    
    \maketitle
    
    

    
    \section{Fast-Fourier Transform (FTT) and Number-Theoretic Transform
(NTT) - FFT, NTT and
LBC}\label{fast-fourier-transform-ftt-and-number-theoretic-transform-ntt---fft-ntt-and-lbc}

    The need to make systems and algorithms faster so as to make them more
practical has been ever constant in the world of software engineering
and cryptography at large, and lattice-based cryptography is not left
out.

Number-theoretic Transform (NTT) was used in ML-KEM
(Module-Lattice-Based Key-Encapsulation Mechanism), a lattice-based
cryptographic algorithm used in establishing a shared secret key between
two parties over a public channel. NTT is an analogue to the beautiful
yet powerful Fast-fourier Transform (FTT).

This article is part one of two-part series titled \textbf{Fast-fourier
Transform (FTT), Number-theoretic Transform (NTT) and Lattice-based
Cryptography}. The goal of this post is to take a deep dive into FFT and
NTT while part two will explore how NTT is used in lattice-based
algorithms like
\href{https://github.com/Godspower-Eze/pqc-ml-kem.rs}{ML-KEM} and
\href{https://csrc.nist.gov/pubs/fips/204/final}{ML-DSA}.

As always, we are going to take a discovery approach, building up slowly
from the familiar to the unfamiliar.

    \subsection{Fast-fourier Transform
(FFT)}\label{fast-fourier-transform-fft}

    Given two polynomials \(A\) and \(B\) of degree 2 : -
\(A = 3 - 4x + x^2\) - \(B = 6 - 5x + x^2\)

    We are asked to multiply them. Easy, right?

\begin{itemize}
\tightlist
\item
  \(C = A\ast B = (3 - 4x + x^2)(6 - 5x + x^2) = 3 * (6 - 5x + x^2) - 4x * (6 - 5x + x^2) + x^2 * (6 - 5x + x^2) = 18 + 39x + 29x^2 - 9x^3 + x^4\)
\end{itemize}

    Below is an implementation of this in Python:

\begin{Shaded}
\begin{Highlighting}[]
\NormalTok{A }\OperatorTok{=}\NormalTok{ [}\DecValTok{3}\NormalTok{, }\OperatorTok{{-}}\DecValTok{4}\NormalTok{, }\DecValTok{1}\NormalTok{] }\CommentTok{\# coefficients of A}

\NormalTok{B }\OperatorTok{=}\NormalTok{ [}\DecValTok{6}\NormalTok{, }\OperatorTok{{-}}\DecValTok{5}\NormalTok{, }\DecValTok{1}\NormalTok{] }\CommentTok{\# coefficients of B}

\NormalTok{d }\OperatorTok{=} \DecValTok{2} \CommentTok{\# degree}

\NormalTok{C }\OperatorTok{=}\NormalTok{ [}\DecValTok{0} \ControlFlowTok{for}\NormalTok{ \_ }\KeywordTok{in} \BuiltInTok{range}\NormalTok{((}\DecValTok{2} \OperatorTok{*}\NormalTok{ d) }\OperatorTok{+} \DecValTok{1}\NormalTok{)]}

\ControlFlowTok{for}\NormalTok{ i }\KeywordTok{in} \BuiltInTok{range}\NormalTok{(d }\OperatorTok{+} \DecValTok{1}\NormalTok{):}
    \ControlFlowTok{for}\NormalTok{ j }\KeywordTok{in} \BuiltInTok{range}\NormalTok{(d }\OperatorTok{+} \DecValTok{1}\NormalTok{):}
\NormalTok{        C[i }\OperatorTok{+}\NormalTok{ j] }\OperatorTok{=}\NormalTok{ A[i] }\OperatorTok{*}\NormalTok{ B[j]}
\end{Highlighting}
\end{Shaded}

    This algorithm takes \(O(d^2)\) time. Is there a faster way to perform
polynomial multiplication?

From the implementation above, you can tell that the polynomial is
represented as a list of its \emph{coefficients}. There are other
representations.

According to the \textbf{uniqueness theorem,} a polynomial of degree at
most \(d\) is uniquely determined by its values at \(d + 1\)
\textbf{distinct} points. That is, a polynomial \(f(x)\) of degree \(d\)
can be uniquely constructed by the points:
\((x_0, f(x_0)), (x_1, f(x_1)),...,(x_d, f(x_d))\). Let's call this
representation \textbf{point representation}.

A more general version of the uniqueness theorem states that a
polynomial of degree at most \(d\) is uniquely determined by its values
at \emph{at least} \(d + 1\) \textbf{distinct} points. That means we
could uniquely represent our polynomial using more than \(d + 1\)
points.

Now, let's represent our polynomials using the point representation,
using the values \(x = -5, -4, -3, -2, -1\). Therefore, our polynomial
can be represented as follows:

{\def\LTcaptype{none} % do not increment counter
\begin{longtable}[]{@{}lllll@{}}
\toprule\noalign{}
\(x\) & \(A(x)\) & \(B(x)\) & \((x, A(x))\) & \((x, B(x))\) \\
\midrule\noalign{}
\endhead
\bottomrule\noalign{}
\endlastfoot
\(-5\) & \(48\) & \(56\) & \((-5, 48)\) & \((-5, 56)\) \\
\(-4\) & \(35\) & \(42\) & \((-4, 35)\) & \((-4, 42)\) \\
\(-3\) & \(24\) & \(30\) & \((-3, 24)\) & \((-3, 30)\) \\
\(-2\) & \(15\) & \(20\) & \((-2, 15)\) & \((-2, 20)\) \\
\(-1\) & \(8\) & \(12\) & \((-1, 8)\) & \((-1, 12)\) \\
\end{longtable}
}

Now, on to the interesting part. It turns out that
\(C(x_i) =A(x_i) \ast B(x_i)\). That is, if you multiply the evaluation
\(A(x_i)\) by the evaluation \(B(x_i)\), you get the evaluation
\(C(x_i)\). See the table below:

{\def\LTcaptype{none} % do not increment counter
\begin{longtable}[]{@{}
  >{\raggedright\arraybackslash}p{(\linewidth - 12\tabcolsep) * \real{0.0571}}
  >{\raggedright\arraybackslash}p{(\linewidth - 12\tabcolsep) * \real{0.0857}}
  >{\raggedright\arraybackslash}p{(\linewidth - 12\tabcolsep) * \real{0.0857}}
  >{\raggedright\arraybackslash}p{(\linewidth - 12\tabcolsep) * \real{0.1571}}
  >{\raggedright\arraybackslash}p{(\linewidth - 12\tabcolsep) * \real{0.1571}}
  >{\raggedright\arraybackslash}p{(\linewidth - 12\tabcolsep) * \real{0.2857}}
  >{\raggedright\arraybackslash}p{(\linewidth - 12\tabcolsep) * \real{0.1714}}@{}}
\toprule\noalign{}
\begin{minipage}[b]{\linewidth}\raggedright
\(x\)
\end{minipage} & \begin{minipage}[b]{\linewidth}\raggedright
\(A(x)\)
\end{minipage} & \begin{minipage}[b]{\linewidth}\raggedright
\(B(x)\)
\end{minipage} & \begin{minipage}[b]{\linewidth}\raggedright
\((x, A(x))\)
\end{minipage} & \begin{minipage}[b]{\linewidth}\raggedright
\((x, B(x))\)
\end{minipage} & \begin{minipage}[b]{\linewidth}\raggedright
\(C(x) = A(x) * B(x)\)
\end{minipage} & \begin{minipage}[b]{\linewidth}\raggedright
\((x, C(x))\)
\end{minipage} \\
\midrule\noalign{}
\endhead
\bottomrule\noalign{}
\endlastfoot
\(-5\) & \(48\) & \(56\) & \((-5, 48)\) & \((-5, 56)\) & \(2688\) &
\((-5, 2688)\) \\
\(-4\) & \(35\) & \(42\) & \((-4, 35)\) & \((-4, 42)\) & \(1470\) &
\((-4, 1470)\) \\
\(-3\) & \(24\) & \(30\) & \((-3, 24)\) & \((-3, 30)\) & \(720\) &
\((-3, 720)\) \\
\(-2\) & \(15\) & \(20\) & \((-2, 15)\) & \((-2, 20)\) & \(300\) &
\((-2, 300)\) \\
\(-1\) & \(8\) & \(12\) & \((-1, 8)\) & \((-1, 12)\) & \(96\) &
\((-1, 96)\) \\
\end{longtable}
}

This means that the points \((x, C(x))\) uniquely represents the
polynomial \(C = 18 + 39x + 29x^2 - 9x^3 + x^4\).

What does this all mean? We just achieved an \(O(d)\) time polynomial
multiplication by using the point representation form of the
polynomials.

Does that mean we just achieved \(O(d)\) for polynomial multiplication?
Not really.

Given two polynomials \(f(x)\) and \(g(x)\) of the form:
\(a_0 + a_1x + a_2x^2 + \cdots + a_{d}x^{d}\), here's the full process:

\begin{enumerate}
\def\labelenumi{\arabic{enumi}.}
\item
  \textbf{Convert from coefficient to point representation}: we pick a
  set of \(x\) values and evaluate the polynomials at those points
  thereby converting the polynomials to the point representation.

  That is, we pick \(x_i = \{x_0, x_1, ... ,x_d \}\). Then, we compute
  \(f(x_i) = \{f(x_0), f(x_1),...,f(x_d)\}\) and
  \(g(x_i) = \{g(x_0), g(x_1),...,g(x_d)\}\).

  There's a way of expressing this in a matrix form:

  \[ V = \begin{bmatrix}x_0^0 & x_0^1 & x_0^2 & \cdots & x_0^{n} \\ x_1^0 & x_1^1 & x_1^2 & \cdots & x_1^{n} \\ x_2^0 & x_2^1 & x_2^2 & \cdots & x_2^{n} \\[0.3em]\vdots & \vdots & \vdots & \ddots & \vdots \\[0.3em]x_n^0 & x^1_{n} & x_{n}^2 & \cdots & x_{n}^{n}\end{bmatrix} = \begin{bmatrix}1 & x_0^1 & x_0^2 & \cdots & x_0^{n} \\ 1 & x_1^1 & x_1^2 & \cdots & x_1^{n} \\ 1 & x_2^1 & x_2^2 & \cdots & x_2^{n} \\[0.3em]\vdots & \vdots & \vdots & \ddots & \vdots \\[0.3em]1 & x^1_{n} & x_{n}^2 & \cdots & x_{n}^{n}\end{bmatrix}\]

  \[a = \begin{bmatrix}a_0 \\ a_1 \\ a_2 \\ \vdots \\ a_n\end{bmatrix}\]

  \[Va = \begin{bmatrix}a_0x_0^0 + a_1x_0^1 + a_2x_0^2 + \cdots + a_nx_0^n\\ a_0x_1^0 + a_1x_1^1 + a_2x_1^2 + \cdots + a_nx_1^n \\ a_0x_2^0 + a_1x_2^1 + a_2x_2^2 + \cdots + a_nx_2^n \\ \vdots \\ a_0x_n^0 + a_1x_n^1 + a_2x_n^2 + \cdots + a_nx_n^n \end{bmatrix} = \begin{bmatrix} f(x_0) \\ f(x_1) \\ f(x_2) \\ \vdots \\ f(x_n) \end{bmatrix}\]

  \[ Va = f(x_i)\]

  where \(n\) is at least \(d\), \(V\) is the matrix of \(x_i\) and
  their powers, \(a\) is a vector of coefficients and \(f(x_i)\) are the
  evaluations. \(V\) is called the \textbf{vandermonde} matrix.

  This shows that the process of evaluating the polynomial \(f(x_i)\) is
  a \emph{matrix-vector multiplication} which is an \(O(n^2)\)
  operation. \(f(x_i)\) can also be written as follows:
  \[f(x_i) = \sum^{n}_{j=0}a_jx^j_i\]

  Keep this in mind, as this formula for expressing \(f(x_i)\) is
  important in understanding FFT later.

  Using \(A(x)\) as an example, choosing \(n = 4\) and
  \(x =\{x_0, x_1, x_2, x_3, x_4\}= \{-5, -4, -3, -2, -1\}\), we have
  the following:

  \[V = \begin{bmatrix}1 & x_0 & x_0^2 & x_0^3 & x_0^4 \\1 & x_1 & x_1^2 & x_1^3 & x_1^4 \\1 & x_2 & x_2^2 & x_2^3 & x_2^4 \\1 & x_3 & x_3^2 & x_3^3 & x_3^4 \\1 & x_4 & x_4^2 & x_4^3 & x_4^4 \end{bmatrix} = \begin{bmatrix}1 & -5 & -5^2 & -5^3 & -5^4 \\ 1 & -4 & -4^2 & -4^3 & -4^4 \\ 1 & -3 & -3^2 & -3^3 & -3^4 \\ 1 & -2 & -2^2 & -2^3 & -2^4 \\ 1 & -1 & -1^2 & -1^3 & -1^4 \end{bmatrix} = \begin{bmatrix}1 & -5 & 25 & -125 & 625 \\1 & -4 & 16 & -64 & 256 \\1 & -3 & 9 & -27 & 81 \\1 & -2 & 4 & -8 & 16 \\1 & -1 & 1 & -1 & 1\end{bmatrix}\]

  \[a = \begin{bmatrix}a_0 \\ a_1 \\ a_2 \\ a_3 \\ a_4\end{bmatrix} = \begin{bmatrix}3 \\ -4 \\ 1 \\ 0 \\ 0\end{bmatrix}\]

  \[Va = A(x_i) = \begin{bmatrix}A(x_0) \\ A(x_1) \\ A(x_2) \\ A(x_3) \\ A(x_4)\end{bmatrix} = \begin{bmatrix}48 \\ 35 \\ 24 \\ 15 \\ 8 \end{bmatrix}\]
\item
  \textbf{Multiply pairwise} (\emph{i.e} \(C(x) = A(x) \ast B(x)\)).
  This runs in \(O(n)\) time.
\item
  \textbf{Convert back from the point represention to the coefficient
  representation}: This is the reverse of step one. We want to find
  \(a\), given \(V\) and \(f(x_i)\). To compute this, we find the
  \textbf{inverse of \(V\)} denoted by \(V^{-1}\) and multiply by
  \(f(x_i)\). That is:

  \[Va = f(x_i)\]

  \[V^{-1}Va = V^{-1}f(x_i)\]

  \[a = V^{-1}f(x_i)\]

  Using \(A(x)\) as an example:

  \[V^{-1} = \begin{bmatrix}1 & -5 & 10 & -10 & 5 \\ 25/12 & -61/6 & 39/2 & -107/6 & 77/12 \\ 35/24 & -41/6 & 49/4 & -59/6 & 71/24 \\5/12 & -11/6 & 3 & -13/6 & 7/12 \\1/24 & -1/6 & 1/4 & -1/6 & 1/24 \end{bmatrix}\]

  \[A(x_i) = \begin{bmatrix}48 \\ 35 \\ 24 \\ 15 \\ 8 \end{bmatrix}\]

  \[a = V^{-1}A(x_i) = \begin{bmatrix}1 & -5 & 10 & -10 & 5 \\ 25/12 & -61/6 & 39/2 & -107/6 & 77/12 \\ 35/24 & -41/6 & 49/4 & -59/6 & 71/24 \\5/12 & -11/6 & 3 & -13/6 & 7/12 \\1/24 & -1/6 & 1/4 & -1/6 & 1/24 \end{bmatrix} \begin{bmatrix}48 \\ 35 \\ 24 \\ 15 \\ 8 \end{bmatrix} = \begin{bmatrix}3 \\ -4 \\ 1 \\ 0 \\ 0\end{bmatrix}\]

  The best version of this algorithm runs in \(O(n^2)\) time.
\end{enumerate}

In total, this process still takes \(O(n^2)\) time.

But, what if there was actually a faster way? A way that allows us to do
polynomial multiplication in less than \(O(n^2)\).

Fast-fourier Transform (FFT) gives us that. FFT allows us to do step one
and step three in \(O(n\log n)\) time; making the whole process
\(O(n \log n)\) time.

To understand FFT, we need to understand the concept of \textbf{complex
numbers} and \textbf{roots of unity}

    \subsubsection{Complex Numbers and Roots of
Unity}\label{complex-numbers-and-roots-of-unity}

    The magic of FFT is based on the beautiful concept of roots of unity but
that in turn is only made possible by the unique properties of complex
numbers. So, here's a deep dive on complex numbers and roots of unity.

Given the polynomial \(P(x) = 1 + x + x^2\), find the roots(\emph{i.e},
the values of \(x\) for which \(P(x) = 0\)) of the polynomial. The
simple answer is the roots does not exist.

Using the quadratic formula \(x = \dfrac{-b \pm \sqrt{b^2 - 4ac}}{2a}\),
let's solve anyways: - \(a = 1\), \(b = 1\) and \(c = 1\) -
\(x = \dfrac{-1 \pm \sqrt{1^2 - 4(1)(1)}}{2(1)} = \dfrac{-1 \pm \sqrt{-3}}{2} = \dfrac{-1}{2} \pm \dfrac{\sqrt{-3}}{2}\)\\
Therefore, the roots of \(P(x)\) are
\(\dfrac{-1}{2} + \dfrac{\sqrt{-3}}{2}\) and
\(\dfrac{-1}{2} - \dfrac{\sqrt{-3}}{2}\).

But, we previously said the polynomial didn't have a root. How come we
got two roots? The answer lies in the fact that from the beginning we
have been dealing with real numbers \(\mathbb{R}\) and the roots
\(\dfrac{-1}{2} + \dfrac{\sqrt{-3}}{2}\) and
\(\dfrac{-1}{2} - \dfrac{\sqrt{-3}}{2}\) are not real numbers.
Therefore, the roots does not exist in real space.

If the roots are not real numbers, then what are they? They are complex
numbers!

Let's further simplify the roots. \(\sqrt{-3}\) can be written
\(\sqrt{3} * \sqrt{-1}\). But, we know that \(\sqrt{-1}\) is impossible
to find. It has a special name, the \emph{imaginary unit} denoted as
\(\mathrm{i}\).

We can rewrite our roots as follows:
\(\dfrac{-1}{2} + \dfrac{\sqrt{3}\mathrm{i}}{2}\) and
\(\dfrac{-1}{2} - \dfrac{\sqrt{3}\mathrm{i}}{2}\). Therefore, a complex
number \(\mathbb{C}\) is a number of the form: \(a + b\mathrm{i}\) where
\(a\) and \(b\) are real numbers and \(\mathrm{i}\) is the imaginary
unit. \(a\) is called the \emph{real part} and \(b\mathrm{i}\) is called
the \emph{imaginary part}. This form of representing a complex number is
called the \emph{rectangular form}. You'll see why soon.

The set of real numbers \(\mathbb{R}\) is a proper subset of the complex
numbers \(\mathbb{C}\) (\emph{i.e}, \(\mathbb{R} \subset \mathbb{C}\)).
Every real number can be written as a complex number where \(b = 0\)
(\emph{e.g} \(5 = 5 + 0\mathrm{i}\)).

One way to describe a real number is as follows: a real number is any
number that can be located on the infinite number line. That means, real
numbers are represented on a number line as shown below.

    \begin{center}
    \adjustimage{max size={0.9\linewidth}{0.9\paperheight}}{ntt_files/ntt_9_0.png}
    \end{center}
    { \hspace*{\fill} \\}
    
    How do we represent a complex number? The number line would have been a
good option but it would only be able to take the real part. We need
diagram that allows us to show both the real part and the imaginary
part. That's where the \textbf{Argand diagram} comes in!

The Argand diagram is two dimensional plane with the real part on the
horinzontal line(x-axis) and the imaginary part on the vertical line
(y-axis).

For example, the complex number \(2 + 3\mathrm{i}\) is represented as
follows:

    \begin{center}
    \adjustimage{max size={0.9\linewidth}{0.9\paperheight}}{ntt_files/ntt_11_0.png}
    \end{center}
    { \hspace*{\fill} \\}
    
    Notice how the dotted lines from \(2\) and \(3\) forms a rectangle.
That's why the form \(a + b\mathrm{i}\) is called the \emph{rectangular
form}.

Imagine we were given the diagram with just the line from the origin as
shown below. Would we be able to identify the complex number?

    \begin{center}
    \adjustimage{max size={0.9\linewidth}{0.9\paperheight}}{ntt_files/ntt_13_0.png}
    \end{center}
    { \hspace*{\fill} \\}
    
    Notice how the angle forms a \emph{right angle triangle}. The diagram
below should help.

    \begin{center}
    \adjustimage{max size={0.9\linewidth}{0.9\paperheight}}{ntt_files/ntt_15_0.png}
    \end{center}
    { \hspace*{\fill} \\}
    
    What do we know about right angle triangles?

According to the \textbf{Pythagoras Theorem}, in a right angled
triangle, the square of the length of the hypotenuse(the side opposite
the right angle) is equal to the sum of the squares of the lengths of
the other two sides(\emph{opposite} and \emph{adjacent}).
Mathematically, if a right-angled triangle has sides \(a\) and \(b\),
and hypotenuse \(c\), then: \(c^2 = a^2 + b^2\).

Having known this; let's adjust out diagram.

    \begin{center}
    \adjustimage{max size={0.9\linewidth}{0.9\paperheight}}{ntt_files/ntt_17_0.png}
    \end{center}
    { \hspace*{\fill} \\}
    
    From the diagram, we can see that we have two unknowns, the hypotenuse
which will be denoted as \(r\) and the angle \(\theta\).

From pythagoras theorem, we know \(r^2 = a^2 + b^2\).
Therefore:\[r^2 = 2^2 + 3^2 = 4 + 9 = 13\] \[r = \sqrt{13}\]

Now we know \(r\), how do we find \(\theta\)? Do you remember the
acronym SOHCAHTOA from high school? This implies:
\[\sin\theta = \dfrac{\mathrm{opposite}}{\mathrm{hypotenuse}} = \dfrac{b}{r},\qquad b = r\sin\theta\]
\[\cos\theta = \dfrac{\mathrm{adjacent}}{\mathrm{hypotenuse}} = \dfrac{a}{r},\qquad a = r\cos\theta\]
\[\tan\theta = \dfrac{\mathrm{opposite}}{\mathrm{adjacent}} = \dfrac{b}{a}, \qquad \theta = \tan^{-1}(\dfrac{b}{a})\]

From this, we can deduce that:
\(a + b\mathrm{i} = r\cos\theta + \mathrm{i}r\sin\theta = r(\cos\theta + \mathrm{i}\sin\theta)\).
This form of representing a complex number is called the \emph{polar
form}.

Now we find \(\theta\) using \(\tan^{-1}(\dfrac{b}{a})\),
\(\quad \theta = \tan^{-1}(\dfrac{3}{2}) = 0.983\) radians.

We have that \(r = \sqrt{13}\) and \(\theta = 0.983\). Given these two
values, we represent the complex number \(2 + 3\mathrm{i}\) as
\(\sqrt{13}(\cos(0.983) + \mathrm{i}\sin(0.983))\). \(0.983\) radians is
\(56.3^\circ\) in degrees and it's better for visualization as shown
below.

The polar representation might not look like the simplest way to express
a complex number but it's going to help us understand roots of unity.

Formally, \(r\) is distance from the point of the complex number to the
origin and it is called the \emph{modulus} or the \emph{absolute value}.
\(\theta\) is called the \emph{argument} of the complex number and it is
the angle that the modulus \(r\) makes from the positive real axis.

    \begin{center}
    \adjustimage{max size={0.9\linewidth}{0.9\paperheight}}{ntt_files/ntt_19_0.png}
    \end{center}
    { \hspace*{\fill} \\}
    
    Let's look at another representation of complex numbers.

Have you ever wondered how calculators computed trignometric functions
like \(\sin(x)\), \(\cos(x)\) and the exponential function \(e^x\)? The
answer lies in approximation using power series.

\(\cos(x)\) and \(\sin(x)\) are approximated using the power series
\(1\ -\frac{x^{2}}{2!} + \frac{x^{4}}{4!} - \frac{x^{6}}{6!} + \frac{x^{8}}{8!} + \cdots\)
and
\(x\ -\frac{x^{3}}{3!} + \frac{x^{5}}{5!} - \frac{x^{7}}{7!} + \frac{x^{9}}{9!} + \cdots\)
respectively. I urge you to verify this using a graph(\emph{i.e} compare
their graphs against each other while extending the series).

Let's substitute \(\theta\) for \(x\):

\[\cos(\theta) = 1\ - \frac{\theta^{2}}{2!} + \frac{\theta^{4}}{4!} - \frac{\theta^{6}}{6!} + \frac{\theta^{8}}{8!} + \cdots\]

\[\sin(\theta) = \theta\ - \frac{\theta^{3}}{3!} + \frac{\theta^{5}}{5!} - \frac{\theta^{7}}{7!} + \frac{\theta^{9}}{9!} + \cdots\]

Now, let's subsitute this approximation into the polar form:

\[\cos\theta + \mathrm{i}\sin\theta = (1\ - \frac{\theta^{2}}{2!} + \frac{\theta^{4}}{4!} - \frac{\theta^{6}}{6!} + \frac{\theta^{8}}{8!} + \cdots) + \mathrm{i}(\theta\ - \frac{\theta^{3}}{3!} + \frac{\theta^{5}}{5!} - \frac{\theta^{7}}{7!} + \frac{\theta^{9}}{9!} + \cdots)\]

\[ \qquad  \qquad  = (1\ - \frac{\theta^{2}}{2!} + \frac{\theta^{4}}{4!} - \frac{\theta^{6}}{6!} + \frac{\theta^{8}}{8!} + \cdots) + (\mathrm{i}\theta\ - \mathrm{i}\frac{\theta^{3}}{3!} + \mathrm{i}\frac{\theta^{5}}{5!} - \mathrm{i}\frac{\theta^{7}}{7!} + \mathrm{i}\frac{\theta^{9}}{9!} + \cdots)\]

\[ \qquad  \qquad  = 1\ + \mathrm{i}\theta - \frac{\theta^{2}}{2!} - \mathrm{i}\frac{\theta^{3}}{3!} + \frac{\theta^{4}}{4!} + \mathrm{i}\frac{\theta^{5}}{5!} - \frac{\theta^{6}}{6!} - \mathrm{i}\frac{\theta^{7}}{7!} + \frac{\theta^{8}}{8!} + \mathrm{i}\frac{\theta^{9}}{9!} + \cdots\]

Let's look at that of the exponential function \(e^x\). The power series
for \(e^x\) as follows:
\[1 + x + \frac{x^2}{2!} + \frac{x^3}{3!} + \frac{x^4}{4!} + \frac{x^5}{5!} + \frac{x^6}{6!} + \frac{x^7}{7!} + \frac{x^8}{8!} + \frac{x^9}{9!} \cdots\]

Let's substitute \(\mathrm{i}\theta\) for \(x\):
\[ e^{\mathrm{i}\theta} = 1 + \mathrm{i}\theta + \frac{(\mathrm{i}\theta)^2}{2!} + \frac{(\mathrm{i}\theta)^3}{3!} + \frac{(\mathrm{i}\theta)^4}{4!} + \frac{(\mathrm{i}\theta)^5}{5!} + \frac{(\mathrm{i}\theta)^6}{6!} + \frac{(\mathrm{i}\theta)^7}{7!} + \frac{(\mathrm{i}\theta)^8}{8!} + \frac{(\mathrm{i}\theta)^9}{9!} + \cdots\]

We will compute the powers of \((\mathrm{i}\theta)\). From indices we
know that \((\mathrm{i}\theta)^n = \theta^n\mathrm{i}^n\). Let's focus
on \(\mathrm{i}^n\): - \(\mathrm{i}^1 = \mathrm{i}\) -
\(\mathrm{i}^2 = \sqrt{-1} \ast \sqrt{-1} = -1\) -
\(\mathrm{i}^3 = \sqrt{-1} \ast \sqrt{-1} \ast \sqrt{-1} = -1 \ast \sqrt{-1} =  -\mathrm{i}\)
-
\(\mathrm{i}^4 = \sqrt{-1} \ast \sqrt{-1} \ast \sqrt{-1} \ast \sqrt{-1} = -1 \ast -1 = 1\)
- \(\mathrm{i}^5 = \mathrm{i}\) - \(\mathrm{i}^6 = -1\) -
\(\mathrm{i}^7 = -\mathrm{i}\) - \(\mathrm{i}^8 = 1\) -
\(\mathrm{i}^9 = \mathrm{i}\)

Do you see the pattern? The values are cyclic in nature; \(\mathrm{i}\)
to \(1\), it repeats and continues like that. We will come back to this
but for now let's go back to computing \(e^{\mathrm{i}\theta}\). It now
becomes:
\[e^{\mathrm{i}\theta} = 1 + \mathrm{i}\theta - \frac{\theta^2}{2!} - \mathrm{i}\frac{\theta^3}{3!} + \frac{\theta^4}{4!} + \mathrm{i}\frac{\theta^5}{5!} - \frac{\theta^6}{6!} - \mathrm{i}\frac{\theta^7}{7!} + \frac{\theta^8}{8!} + \mathrm{i}\frac{\theta^9}{9!} + \cdots\]

From this we can see that
\(e^{\mathrm{i}\theta} = \cos(\theta) + \mathrm{i}\sin(\theta)\).
Therefore, we have our third representation, the \emph{exponential
form}: \(re^{\mathrm{i}\theta}\).

In summary, we have: - rectangular form: \(a + b\mathrm{i}\) - polar
form: \(r(\cos\theta + \mathrm{i}\sin\theta)\) - exponential form:
\(re^{\mathrm{i}\theta}\)

Soon, you will see why exponential and polar form are best suited for
understanding and working with roots of unity.

Now, back to where this it all started: finding the roots of the
polynomial \(P(x) = 1 + x + x^2\). We saw that the roots are
\(\dfrac{-1}{2} + \dfrac{\sqrt{3}\mathrm{i}}{2}\) and
\(\dfrac{-1}{2} - \dfrac{\sqrt{3}\mathrm{i}}{2}\).

Let's represent them using the argand diagram.

    \begin{center}
    \adjustimage{max size={0.9\linewidth}{0.9\paperheight}}{ntt_files/ntt_21_0.png}
    \end{center}
    { \hspace*{\fill} \\}
    
    Let's find the modulus \(r\) and the argument \(\theta\) of these roots:
- For \(\dfrac{-1}{2} + \dfrac{\sqrt{3}\mathrm{i}}{2}\): - finding
\(r\): \(r^2 = a^2 + b^2\), \(\quad a = \dfrac{-1}{2}\),
\(\quad b =  \dfrac{\sqrt{3}}{2}\),
\(\quad r^2 = (\dfrac{-1}{2})^2 + (\dfrac{\sqrt{3}}{2})^2 = \dfrac{1}{4} + \dfrac{3}{4} = \dfrac{4}{4} = 1\),
\(\quad r = 1\) - finding \(\theta\):
\(\theta = \tan^{-1}(\dfrac{b}{a}) = \tan^{-1}(\dfrac{\dfrac{\sqrt{3}}{2}}{\dfrac{-1}{2}}) = \tan^{-1}(\dfrac{\sqrt{3}}{-1}) = \tan^{-1}(-\sqrt{3}) = -60^\circ\)
or \(-\dfrac{\pi}{3}\) - For
\(\dfrac{-1}{2} - \dfrac{\sqrt{3}\mathrm{i}}{2}\): - finding \(r\):
\(r^2 = a^2 + b^2\), \(\quad a = \dfrac{-1}{2}\),
\(\quad b =  \dfrac{-\sqrt{3}}{2}\),
\(\quad r^2 = (\dfrac{-1}{2})^2 + (\dfrac{-\sqrt{3}}{2})^2 = \dfrac{1}{4} + \dfrac{3}{4} = \dfrac{4}{4} = 1\),
\(\quad r = 1\) - finding \(\theta\):
\(\theta = \tan^{-1}(\dfrac{b}{a}) = \tan^{-1}(\dfrac{\dfrac{-\sqrt{3}}{2}}{\dfrac{-1}{2}}) = \tan^{-1}(\dfrac{-\sqrt{3}}{-1}) = \tan^{-1}(\sqrt{3}) = 60^\circ\)
or \(\dfrac{\pi}{3}\)

Let's show our \(r\) and \(\theta\) values to our diagram.

    \begin{center}
    \adjustimage{max size={0.9\linewidth}{0.9\paperheight}}{ntt_files/ntt_23_0.png}
    \end{center}
    { \hspace*{\fill} \\}
    
    This is our argand diagram represented in respect to \(r\) and
\(\theta\). But, it's not exactly correct.

Recall that the argument \(\theta\) is ``the angle that the modulus
makes from the positive real axis''. The key phrase is \emph{from the
positive real axis}.

Let's show this:

    \begin{center}
    \adjustimage{max size={0.9\linewidth}{0.9\paperheight}}{ntt_files/ntt_25_0.png}
    \end{center}
    { \hspace*{\fill} \\}
    
    As you can see, \(\theta\) for
\(\dfrac{-1}{2} + \dfrac{\sqrt{3}\mathrm{i}}{2}\) becomes \(120^\circ\)
or \(\dfrac{2\pi}{3}\) and \(\theta\) for
\(\dfrac{-1}{2} - \dfrac{\sqrt{3}\mathrm{i}}{2}\) becomes \(240^\circ\)
or \(\dfrac{4\pi}{3}\).

In exponential form, this will be written as \(e^{2\pi\mathrm{i}/3}\)
and \(e^{4\pi\mathrm{i}/3}\) respectively as \(r = 1\).

This also helps us see the argument \(\theta\) not just as angle but a
\emph{rotation} from the positive real axis and this is the mental model
we need to understand roots of unity.

Let's do a quick experiment: we are going take the roots individually
and compute their powers up to the third power. Then, we are going to
see what happens in the argand diagram.

Starting with \(e^{2\pi\mathrm{i}/3}\):

\begin{itemize}
\tightlist
\item
  We compute the zero power: \((e^{2\pi\mathrm{i}/3})^0 = e^{0} = 1\)
\end{itemize}

    \begin{center}
    \adjustimage{max size={0.9\linewidth}{0.9\paperheight}}{ntt_files/ntt_27_0.png}
    \end{center}
    { \hspace*{\fill} \\}
    
    \begin{itemize}
\tightlist
\item
  The first power: \((e^{2\pi\mathrm{i}/3})^1 = e^{2\pi\mathrm{i}/3}\)
\end{itemize}

    \begin{center}
    \adjustimage{max size={0.9\linewidth}{0.9\paperheight}}{ntt_files/ntt_29_0.png}
    \end{center}
    { \hspace*{\fill} \\}
    
    \begin{itemize}
\tightlist
\item
  The second power \((e^{2\pi\mathrm{i}/3})^2 = e^{4\pi\mathrm{i}/3}\)
\end{itemize}

    \begin{center}
    \adjustimage{max size={0.9\linewidth}{0.9\paperheight}}{ntt_files/ntt_31_0.png}
    \end{center}
    { \hspace*{\fill} \\}
    
    \begin{itemize}
\tightlist
\item
  The third power \((e^{2\pi\mathrm{i}/3})^3 = e^{2\pi\mathrm{i}}\)
\end{itemize}

    \begin{center}
    \adjustimage{max size={0.9\linewidth}{0.9\paperheight}}{ntt_files/ntt_33_0.png}
    \end{center}
    { \hspace*{\fill} \\}
    
    \begin{center}\rule{0.5\linewidth}{0.5pt}\end{center}

    Then, for \(e^{4\pi\mathrm{i}/3}\):

\begin{itemize}
\tightlist
\item
  We compute the zero power: \((e^{4\pi\mathrm{i}/3})^0 = e^0 = 1\)
\end{itemize}

    \begin{center}
    \adjustimage{max size={0.9\linewidth}{0.9\paperheight}}{ntt_files/ntt_36_0.png}
    \end{center}
    { \hspace*{\fill} \\}
    
    \begin{itemize}
\tightlist
\item
  The first power: \((e^{4\pi\mathrm{i}/3})^1 = e^{4\pi\mathrm{i}/3}\)
\end{itemize}

    \begin{center}
    \adjustimage{max size={0.9\linewidth}{0.9\paperheight}}{ntt_files/ntt_38_0.png}
    \end{center}
    { \hspace*{\fill} \\}
    
    \begin{itemize}
\tightlist
\item
  The second power: \((e^{4\pi\mathrm{i}/3})^2 = e^{8\pi\mathrm{i}/3}\)
\end{itemize}

    \begin{center}
    \adjustimage{max size={0.9\linewidth}{0.9\paperheight}}{ntt_files/ntt_40_0.png}
    \end{center}
    { \hspace*{\fill} \\}
    
    \begin{itemize}
\tightlist
\item
  The third power: \((e^{4\pi\mathrm{i}/3})^3 = e^{4\pi\mathrm{i}}\)
\end{itemize}

    \begin{center}
    \adjustimage{max size={0.9\linewidth}{0.9\paperheight}}{ntt_files/ntt_42_0.png}
    \end{center}
    { \hspace*{\fill} \\}
    
    \begin{center}\rule{0.5\linewidth}{0.5pt}\end{center}

    What do you notice?

The powers cycles back to \(1\) at the third power. If we continued this
cycle, you would notice that it doesn't just cycle back to \(1\) at the
third power but every power that's a multiple of \(3\). I urge you to
try this out.

Let's try this again for another polynomial.

Given the polynomial \(G(x) = x^4 - 1\), let's the find the roots.

We get that:

\begin{itemize}
\tightlist
\item
  \(G(x) = x^4 - 1 = (x^2 - 1)(x^2 + 1) = (x - 1)(x + 1)(x^2 + 1)\)
\end{itemize}

Setting them to zero, we have that: - \((x - 1) = 0\), \(x = 1\) -
\((x + 1) = 0\), \(x = -1\) - \((x^2 + 1) = 0\), \(x^2 = -1\),
\(x = \pm\sqrt{-1} = \pm\mathrm{i}\)

Therefore, the roots of \(G(x)\) are \(1\), \(-1\), \(-\mathrm{i}\) and
\(\mathrm{i}\)

Let's show them on the argand diagram.

    \begin{center}
    \adjustimage{max size={0.9\linewidth}{0.9\paperheight}}{ntt_files/ntt_45_0.png}
    \end{center}
    { \hspace*{\fill} \\}
    
    \begin{center}\rule{0.5\linewidth}{0.5pt}\end{center}

    Let's compute the modulus \(r\) and argument \(\theta\) for each of
them.

\begin{itemize}
\tightlist
\item
  For \(1\):

  \begin{itemize}
  \tightlist
  \item
    finding \(r\): \(r^2 = a^2 + b^2\), \(\quad a = 1\),
    \(\quad b = 0\), \(\quad r^2 = 1^2 + 0^2 = 1\), \(\quad r = 1\)
  \item
    finding \(\theta\):
    \(\theta = \tan^{-1}(\dfrac{b}{a}) = \tan^{-1}(\dfrac{0}{1}) = \tan^{-1}(0) = 0 = 0^\circ\)
    or \(0\pi\)
  \end{itemize}
\item
  For \(-1\):

  \begin{itemize}
  \tightlist
  \item
    finding \(r\): \(r^2 = a^2 + b^2\), \(\quad a = -1\),
    \(\quad b = 0\), \(\quad r^2 = -1^2 + 0^2 = 1\), \(\quad r = 1\)
  \item
    finding \(\theta\):
    \(\theta = \tan^{-1}(\dfrac{b}{a}) = \tan^{-1}(\dfrac{0}{-1}) = \tan^{-1}(0) = 0 = 0^\circ\).
    \(0^\circ\) from the positive x-axis is \(180^\circ\) so
    \(\theta = 180^\circ\) or \(\pi\)
  \end{itemize}
\item
  For \(\mathrm{i}\):

  \begin{itemize}
  \tightlist
  \item
    finding \(r\): \(r^2 = a^2 + b^2\), \(\quad a = 0\),
    \(\quad b = 1\), \(\quad r^2 = 0^2 + 1^2 = 1\), \(\quad r = 1\)
  \item
    finding \(\theta\):
    \(\theta = \tan^{-1}(\dfrac{b}{a}) = \tan^{-1}(\dfrac{1}{0}) = 90^\circ\)
    or \(\dfrac{\pi}{2}\). This looks like we are dividing 1 by 0 but
    trust me that's not the case. I will leave you to verify this
    yourself.
  \end{itemize}
\item
  For \(-\mathrm{i}\):

  \begin{itemize}
  \tightlist
  \item
    finding \(r\): \(r^2 = a^2 + b^2\), \(\quad a = 0\),
    \(\quad b = -1\), \(\quad r^2 = 0^2 + (-1)^2 = 1\), \(\quad r = 1\)
  \item
    finding \(\theta\):
    \(\theta = \tan^{-1}(\dfrac{b}{a}) = \tan^{-1}(\dfrac{-1}{0}) = \tan^{-1}(0) = -90^\circ\).
    \(-90^\circ\) from the positive x-axis is \(270\) so
    \(\theta = 270^\circ\) or \(\dfrac{3\pi}{2}\)
  \end{itemize}
\end{itemize}

Converting them to exponential form, being that \(r = 1\) for all roots,
we have that: - \(1 = e^{0\mathrm{i}} = e^0\) -
\(-1 = e^{\pi\mathrm{i}}\) - \(\mathrm{i} = e^{\pi\mathrm{i}/{2}}\) -
\(-\mathrm{i} = e^{3\pi\mathrm{i}/{2}}\)

Now, let's compute the powers of each root up to the fourth power.

Starting with \(e^0\):

\begin{itemize}
\tightlist
\item
  We compute the zero power \((e^0)^0 = e^0\). Notice that all the
  powers are equal as
  \((e^0)^0 = (e^0)^1 = (e^0)^2 = (e^0)^3 = (e^0)^4 = e^0\) so we have
  just one argand diagram showing the point without rotation.
\end{itemize}

    \begin{center}
    \adjustimage{max size={0.9\linewidth}{0.9\paperheight}}{ntt_files/ntt_48_0.png}
    \end{center}
    { \hspace*{\fill} \\}
    
    \begin{center}\rule{0.5\linewidth}{0.5pt}\end{center}

    For \(e^{\pi\mathrm{i}}\): - We compute the zero power:
\((e^{\pi\mathrm{i}})^0 = e^0\)

    \begin{center}
    \adjustimage{max size={0.9\linewidth}{0.9\paperheight}}{ntt_files/ntt_51_0.png}
    \end{center}
    { \hspace*{\fill} \\}
    
    \begin{itemize}
\tightlist
\item
  The first power: \((e^{\pi\mathrm{i}})^1 = e^{\pi\mathrm{i}}\)
\end{itemize}

    \begin{center}
    \adjustimage{max size={0.9\linewidth}{0.9\paperheight}}{ntt_files/ntt_53_0.png}
    \end{center}
    { \hspace*{\fill} \\}
    
    \begin{itemize}
\tightlist
\item
  The second power: \((e^{\pi\mathrm{i}})^2 = e^{2\pi\mathrm{i}}\)
\end{itemize}

    \begin{center}
    \adjustimage{max size={0.9\linewidth}{0.9\paperheight}}{ntt_files/ntt_55_0.png}
    \end{center}
    { \hspace*{\fill} \\}
    
    \begin{itemize}
\tightlist
\item
  The third power: \((e^{\pi\mathrm{i}})^3 = e^{3\pi\mathrm{i}}\)
\end{itemize}

    \begin{center}
    \adjustimage{max size={0.9\linewidth}{0.9\paperheight}}{ntt_files/ntt_57_0.png}
    \end{center}
    { \hspace*{\fill} \\}
    
    \begin{itemize}
\tightlist
\item
  The fourth power: \((e^{\pi\mathrm{i}})^4 = e^{4\pi\mathrm{i}}\)
\end{itemize}

    \begin{center}
    \adjustimage{max size={0.9\linewidth}{0.9\paperheight}}{ntt_files/ntt_59_0.png}
    \end{center}
    { \hspace*{\fill} \\}
    
    \begin{center}\rule{0.5\linewidth}{0.5pt}\end{center}

    For \(e^{\pi\mathrm{i}/{2}}\):

\begin{itemize}
\tightlist
\item
  We compute the zero power: \((e^{\pi\mathrm{i}/{2}})^0 = e^0\)
\end{itemize}

    \begin{center}
    \adjustimage{max size={0.9\linewidth}{0.9\paperheight}}{ntt_files/ntt_62_0.png}
    \end{center}
    { \hspace*{\fill} \\}
    
    \begin{itemize}
\tightlist
\item
  The first power: \((e^{\pi\mathrm{i}/{2}})^1 = e^{\pi\mathrm{i}/{2}}\)
\end{itemize}

    \begin{center}
    \adjustimage{max size={0.9\linewidth}{0.9\paperheight}}{ntt_files/ntt_64_0.png}
    \end{center}
    { \hspace*{\fill} \\}
    
    \begin{itemize}
\tightlist
\item
  The second power: \((e^{\pi\mathrm{i}/{2}})^2 = e^{\pi\mathrm{i}}\)
\end{itemize}

    \begin{center}
    \adjustimage{max size={0.9\linewidth}{0.9\paperheight}}{ntt_files/ntt_66_0.png}
    \end{center}
    { \hspace*{\fill} \\}
    
    \begin{itemize}
\tightlist
\item
  The third power: \((e^{\pi\mathrm{i}/{2}})^3 = e^{3\pi\mathrm{i}/2}\)
\end{itemize}

    \begin{center}
    \adjustimage{max size={0.9\linewidth}{0.9\paperheight}}{ntt_files/ntt_68_0.png}
    \end{center}
    { \hspace*{\fill} \\}
    
    \begin{itemize}
\tightlist
\item
  The fourth power: \((e^{\pi\mathrm{i}/{2}})^4 = e^{2\pi\mathrm{i}}\)
\end{itemize}

    \begin{center}
    \adjustimage{max size={0.9\linewidth}{0.9\paperheight}}{ntt_files/ntt_70_0.png}
    \end{center}
    { \hspace*{\fill} \\}
    
    \begin{center}\rule{0.5\linewidth}{0.5pt}\end{center}

    For \(e^{3\pi\mathrm{i}/{2}}\):

\begin{itemize}
\tightlist
\item
  We compute the zero power: \((e^{3\pi\mathrm{i}/{2}})^0 = e^0\)
\end{itemize}

    \begin{center}
    \adjustimage{max size={0.9\linewidth}{0.9\paperheight}}{ntt_files/ntt_73_0.png}
    \end{center}
    { \hspace*{\fill} \\}
    
    \begin{itemize}
\tightlist
\item
  The first power:
  \((e^{3\pi\mathrm{i}/{2}})^1 = e^{3\pi\mathrm{i}/{2}}\)
\end{itemize}

    \begin{center}
    \adjustimage{max size={0.9\linewidth}{0.9\paperheight}}{ntt_files/ntt_75_0.png}
    \end{center}
    { \hspace*{\fill} \\}
    
    \begin{itemize}
\tightlist
\item
  The second power: \((e^{3\pi\mathrm{i}/{2}})^2 = e^{3\pi\mathrm{i}}\)
\end{itemize}

    \begin{center}
    \adjustimage{max size={0.9\linewidth}{0.9\paperheight}}{ntt_files/ntt_77_0.png}
    \end{center}
    { \hspace*{\fill} \\}
    
    \begin{itemize}
\tightlist
\item
  The third power: \((e^{3\pi\mathrm{i}/{2}})^3 = e^{9\pi\mathrm{i}/2}\)
\end{itemize}

    \begin{center}
    \adjustimage{max size={0.9\linewidth}{0.9\paperheight}}{ntt_files/ntt_79_0.png}
    \end{center}
    { \hspace*{\fill} \\}
    
    \begin{itemize}
\tightlist
\item
  The fourth power: \((e^{3\pi\mathrm{i}/{2}})^4 = e^{6\pi\mathrm{i}}\)
\end{itemize}

    \begin{center}
    \adjustimage{max size={0.9\linewidth}{0.9\paperheight}}{ntt_files/ntt_81_0.png}
    \end{center}
    { \hspace*{\fill} \\}
    
    \begin{center}\rule{0.5\linewidth}{0.5pt}\end{center}

    What do we notice?

\begin{itemize}
\tightlist
\item
  For \(e^0\):

  \begin{itemize}
  \tightlist
  \item
    They are no rotations. It just remains at \(1\)
  \end{itemize}
\item
  For \(e^{\pi\mathrm{i}}\):

  \begin{itemize}
  \tightlist
  \item
    It rotates between \(1\) and \(-1\) and it's at point \(1\) at the
    powers divisible by \(2\)(\(2\) and \(4\)).
  \end{itemize}
\item
  For \(e^{\pi\mathrm{i}/{2}}\):

  \begin{itemize}
  \tightlist
  \item
    It rotates across all roots and meets \(1\) at the fourth power.
  \end{itemize}
\item
  For \(e^{3\pi\mathrm{i}/{2}}\):

  \begin{itemize}
  \tightlist
  \item
    Same as \(e^{\pi\mathrm{i}/{2}}\)
  \end{itemize}
\end{itemize}

Finally, what are roots of unity and why does all this matter?

The \(n\)th roots are unity are simply values of \(x\) that satisfy the
function \(x^n = 1\) for an arbitrary \(n\) value.

Having known this, notice how the roots \(1\), \(-1\), \(\mathrm{i}\)
and \(-\mathrm{i}\) all satisfy function \(x^4 = 1\). That is:

\begin{itemize}
\tightlist
\item
  \(1^4 = 1\) and \((e^0)^4 = 1\)
\item
  \(-1^4 = 1\) and \((e^{\pi\mathrm{i}})^4 = 1\)
\item
  \(\mathrm{i}^4 = 1\) and
  \((e^{\pi\mathrm{i}/{2}})^4 = e^{2\pi\mathrm{i}} = 1\)
\item
  \((-\mathrm{i})^4 = 1\) and
  \((e^{3\pi\mathrm{i}/{2}})^4 = e^{6\pi\mathrm{i}} = 1\)
\end{itemize}

And, how \(x^4 - 1 = 0\) is equals \(x^4 = 1\).

These are \(4\)th roots of unity.

But then, what about the polynomial \(P(x) = 1 + x + x^2\), does that
hold for it too? Let's find out.

Given the polynomial \(x^3 - 1\), let's find the roots.

\begin{itemize}
\tightlist
\item
  \((x^3 - 1) = (x - 1)(1 + x + x^2) = (x - 1)P(x)\)
\end{itemize}

Setting them to zero, we have that: - \(x - 1 = 0\), \(x = 1\) -
\(P(x) = 0\), we already know the roots are
\(\dfrac{-1}{2} + \dfrac{\sqrt{3}\mathrm{i}}{2}\) and
\(\dfrac{-1}{2} - \dfrac{\sqrt{3}\mathrm{i}}{2}\).

Therefore, \(1\), \(\dfrac{-1}{2} + \dfrac{\sqrt{3}\mathrm{i}}{2}\) and
\(\dfrac{-1}{2} - \dfrac{\sqrt{3}\mathrm{i}}{2}\) are the \(3\)rd roots
of unity. That is:

\begin{itemize}
\tightlist
\item
  \(1^3 = 1\) and \((e^0)^3 = 1\)
\item
  \((\dfrac{-1}{2} + \dfrac{\sqrt{3}\mathrm{i}}{2})^3 = 1\) and
  \((e^{2\pi\mathrm{i}/3})^3 = e^{2\pi\mathrm{i}} = 1\)
\item
  \((\dfrac{-1}{2} - \dfrac{\sqrt{3}\mathrm{i}}{2})^3 = 1\) and
  \((e^{4\pi\mathrm{i}/3})^3 = e^{4\pi\mathrm{i}} = 1\)
\end{itemize}

Let's extract more insights from our argand diagrams.

Notice how in the \(3\)rd roots of unity, \(e^{2\pi\mathrm{i}/3}\) and
\(e^{4\pi\mathrm{i}/3}\) generates every other root in the set and in
the \(4\)th roots of unity \(e^{\pi\mathrm{i}/{2}}\) and
\(e^{3\pi\mathrm{i}/{2}}\) also generates every other other roots in the
set. For example, - For \(e^{2\pi\mathrm{i}/3}\) in the \(3\)rd roots of
unity: - \((e^{2\pi\mathrm{i}/3})^ 0 = 1\) -
\((e^{2\pi\mathrm{i}/3})^ 1 = e^{2\pi\mathrm{i}/3}\) -
\((e^{2\pi\mathrm{i}/3})^ 2 = e^{4\pi\mathrm{i}/3}\) - For
\(e^{3\pi\mathrm{i}/{2}}\) in the \(4\)th roots of unity: -
\((e^{3\pi\mathrm{i}/{2}})^ 0 = 1\) -
\((e^{3\pi\mathrm{i}/{2}})^ 1 = e^{3\pi\mathrm{i}/{2}}\) -
\((e^{3\pi\mathrm{i}/{2}})^ 2 = e^{3\pi\mathrm{i}} = e^{\pi\mathrm{i}}\)
-
\((e^{3\pi\mathrm{i}/{2}})^ 3 = e^{9\pi\mathrm{i}/{2}} = e^{\pi\mathrm{i}/{2}}\)

When a root generates all other roots in the set, it is called a
\textbf{primitive root of unity}. This is evident in the diagrams as
\(e^{\pi\mathrm{i}/{2}}\) and \(e^{3\pi\mathrm{i}/{2}}\) hits every root
of unity as it rotates around the circle.

There's something that has been recurring ever since we started talking
about roots of polynomials: the modulus \(r\) has always been \(1\).
This is not a concidence.

If you interpret \(r\) as the radius of a circle being the \textbf{unit
circle}(i.e a circle of radius \(1\)) then every root of unity is a
point on the unit circle hence the name \emph{roots of unity}.

And, there's a general formular for representing every \(n\)th root of
unity: \[e^{{2\pi\mathrm{i} \ast k}/n}\] where \(n\) is an arbitrary
value and \(k = 0, 1, \cdots, n - 1\).

For example:

\begin{itemize}
\tightlist
\item
  The \(5\)th roots of unity are
  \(\{e^{{2\pi\mathrm{i} \ast 0}/5}, e^{{2\pi\mathrm{i}\ast 1}/5}, e^{{2\pi\mathrm{i} \ast 2}/5}, e^{{2\pi\mathrm{i} \ast 3}/5}, e^{{2\pi\mathrm{i} \ast 4}/5}\} = \{e^0, e^{{2\pi\mathrm{i}}/5}, e^{{4\pi\mathrm{i}}/5}, e^{{6\pi\mathrm{i}}/5}, e^{{8\pi\mathrm{i}}/5}\}\)
\end{itemize}

    \begin{center}
    \adjustimage{max size={0.9\linewidth}{0.9\paperheight}}{ntt_files/ntt_84_0.png}
    \end{center}
    { \hspace*{\fill} \\}
    
    \begin{itemize}
\tightlist
\item
  The \(8\)th roots of unity are \$\{ e\^{}\{\{2\pi\mathrm{i}
  \ast 0\}/8\}, e\^{}\{\{2\pi\mathrm{i} \ast 1\}/8\},
  e\^{}\{\{2\pi\mathrm{i} \ast 2\}/8\}, e\^{}\{\{2\pi\mathrm{i}
  \ast 3\}/8\}, e\^{}\{\{2\pi\mathrm{i} \ast 4\}/8\},
  e\^{}\{\{2\pi\mathrm{i} \ast 5\}/8\}, e\^{}\{\{2\pi\mathrm{i}
  \ast 6\}/8\}, e\^{}\{\{2\pi\mathrm{i} \ast 7\}/8\} \} = \{ e\^{}\{0\},
  e\^{}\{\{\pi\mathrm{i}\}/4\}, e\^{}\{\{\pi\mathrm{i}\}/2\},
  e\^{}\{\{3\pi\mathrm{i}\}/4\}, e\^{}\{\pi\mathrm{i}\},
  e\^{}\{\{5\pi\mathrm{i}\}/4\}, e\^{}\{\{3\pi\mathrm{i}\}/2\},
  e\^{}\{\{7\pi\mathrm{i}\}/4\} \} \$
\end{itemize}

    \begin{center}
    \adjustimage{max size={0.9\linewidth}{0.9\paperheight}}{ntt_files/ntt_86_0.png}
    \end{center}
    { \hspace*{\fill} \\}
    
    \begin{center}\rule{0.5\linewidth}{0.5pt}\end{center}

    A common way of representing roots of unity is in terms of the primitive
root of unity. Recall that a primitive root of unity generates all other
roots of unity.

Given a primitive root of unity \(\omega_n\), the roots of unity are
\(\{\omega_n^{0}, \omega_n^{1}, ... , \omega_n^{n - 1}\}\) where \(n\)
is an arbitrary positive number.

For example, from the \(8\)th roots of unity, we pick the primitive root
\(\omega_8 = e^{{\pi\mathrm{i}}/4}\). That \(8\)th roots of unity are as
follows:
\[\{\omega_8^0, \omega_8^1, \omega_8^2, \omega_8^3, \omega_8^4, \omega_8^5, \omega_8^6, \omega_8^7\}\]

    \begin{center}\rule{0.5\linewidth}{0.5pt}\end{center}

    Another important property of roots of unity is that every root has a
\emph{complex conjugate} and that complex conjugate is the point
symmetric to it on the unit circle.

The complex conjugate of a complex number \(a + b\mathrm{i}\) is
\(a - b\mathrm{i}\) and that of \(a - b\mathrm{i}\) is
\(a + b\mathrm{i}\). But then, when these are on the unit circle, the
complex conjugate of any root is simply it's vertical reflexion.

For example, from the \(8\)th roots of unity above, Here are the complex
conjugates: - \(e^{3\pi\mathrm{i}/4}\) is \(e^{5\pi\mathrm{i}/4}\) -
\(e^{\pi\mathrm{i}/4}\) is \(e^{7\pi\mathrm{i}/4}\) -
\(e^{\pi\mathrm{i}/2}\) is \(e^{3\pi\mathrm{i}/2}\)

Here are the \(8\)th roots of unity in rectangular form for
clarification:
\(\{ e^{0}, e^{{\pi\mathrm{i}}/4}, e^{{\pi\mathrm{i}}/2}, e^{{3\pi\mathrm{i}}/4}, e^{\pi\mathrm{i}}, e^{{5\pi\mathrm{i}}/4}, e^{{3\pi\mathrm{i}}/2}, e^{{7\pi\mathrm{i}}/4} \} = \{ 1, \dfrac{\sqrt{2}}{2} + \mathrm{i}\dfrac{\sqrt{2}}{2}, \mathrm{i}, -\dfrac{\sqrt{2}}{2} + \mathrm{i}\dfrac{\sqrt{2}}{2}, -1, -\dfrac{\sqrt{2}}{2} - \mathrm{i}\dfrac{\sqrt{2}}{2}, - \mathrm{i}, \dfrac{\sqrt{2}}{2} - \mathrm{i}\dfrac{\sqrt{2}}{2} \}\)

Interestingly, the complex conjugate of any root \(e^{a\mathrm{i}}\) is
\(e^{-a\mathrm{i}}\). That is:

\begin{itemize}
\tightlist
\item
  \(e^{5\pi\mathrm{i}/4} = e^{-5\pi\mathrm{i}/4} = e^{3\pi\mathrm{i}/4}\)
\item
  \(e^{7\pi\mathrm{i}/4} = e^{-7\pi\mathrm{i}/4} = e^{\pi\mathrm{i}/4}\)
\item
  \(e^{3\pi\mathrm{i}/2} = e^{-3\pi\mathrm{i}/2} = e^{\pi\mathrm{i}/2}\)
\end{itemize}

A good mental model for complex conjugates is anticlockwise vs clockwise
rotations. For example \(e^{3\pi\mathrm{i}/4}\) is the angle \(3\pi/4\)
in the anticlockwise direction while \(e^{5\pi\mathrm{i}/4}\) is the
same angle in clockwise direction.

An even better way to interprete this is \emph{performing a rotation and
undoing that rotation}.

For example, in the diagram below, we have a value \(2\).

    \begin{center}
    \adjustimage{max size={0.9\linewidth}{0.9\paperheight}}{ntt_files/ntt_91_0.png}
    \end{center}
    { \hspace*{\fill} \\}
    
    Let's multiply 2 by \(e^{{3\pi\mathrm{i}}/4}\). We have
\(2e^{{3\pi\mathrm{i}}/4}\).

    \begin{center}
    \adjustimage{max size={0.9\linewidth}{0.9\paperheight}}{ntt_files/ntt_93_0.png}
    \end{center}
    { \hspace*{\fill} \\}
    
    Now, let's multiply \(2e^{{3\pi\mathrm{i}}/4}\) by the complex conjugate
of \(e^{{3\pi\mathrm{i}}/4}\), \(e^{5\pi\mathrm{i}/4}\). That is
\(2e^{{3\pi\mathrm{i}}/4} * e^{5\pi\mathrm{i}/4} = 2e^{{2\pi\mathrm{i}}}\).

    \begin{center}
    \adjustimage{max size={0.9\linewidth}{0.9\paperheight}}{ntt_files/ntt_95_0.png}
    \end{center}
    { \hspace*{\fill} \\}
    
    As you can see in the diagrams \(e^{{3\pi\mathrm{i}}/4}\) rotated \(2\)
by \({3\pi}/4\) and \(e^{{5\pi\mathrm{i}}/4}\) undid that rotation
taking \(2e^{{3\pi\mathrm{i}}/4}\) back to \(2\).

This is a key concept in next our topic: \textbf{Discrete Fourier
Transform} as we get closer to understanding FFT.

    \subsubsection{Discrete Fourier Transform
(DFT)}\label{discrete-fourier-transform-dft}

    Let's go back to where it all started: \emph{Polynomial multiplication}.

Recall, we chose a set of \(x\) values, \(x_i = \{x_0, x_1,..., x_n\}\)
and evaluated \(f(x_i)\) to get the set of points
\((x_0, f(x_0)), (x_1, f(x_1)),...,(x_d, f(x_d))\) and we called it the
\textbf{point representation}.

Now, let's use roots of unit as those set \(x\) values; our evaluation
points. Given a primitive root of unity \(\omega_n\), we have our roots
of unit
\(x_i = \{\omega_n^0, \omega_n^1, \omega_n^2, ..., \omega_n^{n - 1} \}\).

With this, we have the following:

\[ V = \begin{bmatrix} \omega_n^{0*0} & \omega_n^{0*1} & {\omega_n}^{0*2} & \cdots & \omega_n^{0*(n - 1)} \\ \omega_n^{1*0} & \omega_n^{1*1} & \omega_n^{1*2} & \cdots & \omega_n^{1*(n - 1)} \\ \omega_n^{2*0} & \omega_n^{2*1} & \omega_n^{2*2} & \cdots & \omega_n^{2*(n - 1)} \\[0.3em]\vdots & \vdots & \vdots & \ddots & \vdots \\[0.3em]\omega_n^{(n - 1)*0} & \omega_n^{(n - 1)*1} & \omega_n^{(n - 1)*2} & \cdots & \omega_n^{(n - 1)(n - 1)} \end{bmatrix} = \begin{bmatrix}\omega_n^0 & \omega_n^0 & \omega_n^0 & \cdots & \omega_n^0 \\ \omega_n^0 & \omega_n^1 & \omega_n^2 & \cdots & \omega_n^{n - 1} \\ \omega_n^0 & \omega_n^2 & \omega_n^4 & \cdots & \omega_n^{2n - 2} \\[0.3em]\vdots & \vdots & \vdots & \ddots & \vdots \\[0.3em] \omega_n^0 & \omega_n^{n - 1} & \omega_n^{2n - 2} & \cdots & \omega_n^{n^2 - 2n + 1} \end{bmatrix}\]

\[a = \begin{bmatrix}a_0 \\ a_1 \\ a_2 \\ \vdots \\ a_{n - 1}\end{bmatrix}\]

\[Va = \begin{bmatrix}a_0\omega_n^0 + a_1\omega_n^0 + a_2\omega_n^0 + \cdots + a_{n - 1}w_n^0 \\ a_0\omega_n^0 + a_1\omega_n^1 + a_2\omega_n^2 + \cdots + a_{n - 1}\omega_n^{n - 1} \\ a_0\omega_n^0 + a_1\omega_n^2 + a_2\omega_n^4 + \cdots + a_{n - 1}\omega_n^{2n - 2} \\ \vdots \\ a_0\omega_n^0 + a_1\omega_n^{n - 1} + a_2\omega_n^{2n - 2} + \cdots + a_{n - 1}\omega_n^{n^2 - 2n + 1} \end{bmatrix} = \begin{bmatrix} f(\omega_n^0) \\ f(\omega_n^1) \\ f(\omega_n^2) \\ \vdots \\ f(\omega_n^{n - 1}) \end{bmatrix}\]

\[ Va = f(x_i)\]

    Breifly, what is DFT? DFT is an algorithm that turns a signal from the
time domain into the frequency domain. It used in Signal analysis, Image
Compression and a lot more.

A good way to look at DFT is an algorithm that seeks to get more
information from individual values that were created out of a
combination of different measurements of values. For instance, given an
chord(a chord is a set of notes played together), DFT can be used to get
the frequency of the individual notes that makes up the chord.

Why do we care about DFT? The DFT algorithm is equalivent to polynomial
evaluation using roots of unity as we shown above. It runs in \(O(n^2)\)
time and \textbf{FFT is an optimization of DFT and it run in
\(O(n\log(n))\) time}.

We are going to see how DFT is used and after that we are going to see
how FFT is derived from the construction of DFT.

Let's see an example. Using the \(4\)th roots of unity, let's compute
\(C(x) = A(x) * B(x)\) where \(A(x) = 1 + 2x\), \(B(x) = 3 + 4x\) and
\(C(x) = 3 + 10x + 8x^2\).

\begin{itemize}
\tightlist
\item
  \textbf{Step 1}: Pick a primitive root of unity and define the the
  vandermonde matrix \(V\)
\end{itemize}

    \[\omega_4 = e^{{3\pi\mathrm{i}}/2}\]

\[\omega_4 = \{ \omega_4^{0}, \omega_4^{1}, \omega_4^{2}, \omega_4^{3} \} = \{ 1, e^{{3\pi\mathrm{i}}/2}, e^{3\pi\mathrm{i}}, e^{{9\pi\mathrm{i}}/2} \}\]

\[ V = 
     \begin{bmatrix}
     \omega_4^{0*0} & \omega_4^{0*1} & {\omega_4}^{0*2} & \omega_4^{0*3}
     \\ \omega_4^{1*0} & \omega_4^{1*1} & {\omega_4}^{1*2} & \omega_4^{1*3} 
     \\ \omega_4^{2*0} & \omega_4^{2*1} & {\omega_4}^{2*2} & \omega_4^{2*3} 
     \\  \omega_4^{3*0} & \omega_4^{3*1} & {\omega_4}^{3*2} & \omega_4^{3*3}
     \end{bmatrix} = 
     \begin{bmatrix} \omega_4^{0} & \omega_4^{0} & {\omega_4}^{0} & \omega_4^{0}
     \\ \omega_4^{0} & \omega_4^{1} & {\omega_4}^{2} & \omega_4^{3}
     \\ \omega_4^{0} & \omega_4^{2} & {\omega_4}^{4} & \omega_4^{6}
     \\  \omega_4^{0} & \omega_4^{3} & {\omega_4}^{6} & \omega_4^{9}
     \end{bmatrix} = 
     \begin{bmatrix} 1 & 1 & 1 & 1 
     \\ 1 & e^{{3\pi\mathrm{i}}/2} & e^{3\pi\mathrm{i}} & e^{{9\pi\mathrm{i}}/2}
     \\ 1 & e^{3\pi\mathrm{i}} & e^{6\pi\mathrm{i}} & e^{9\pi\mathrm{i}}
     \\ 1 & e^{{9\pi\mathrm{i}}/2} & e^{9\pi\mathrm{i}} & e^{{27\pi\mathrm{i}}/2}
     \end{bmatrix}\]

    \begin{itemize}
\item
  \textbf{Step 2}: Compute the DFT of \(A(x)\) and \(B(x)\),
  \(\mathrm{DFT}(A)\) and \(\mathrm{DFT}(B)\) respectively.

  The formula for computing the DFT is
  \(X[k] = \sum_{n=0}^{N-1} x[n]\, \omega_N^{kn}\) where
  \(\omega_N = e^{-2\pi i / N}\) and \(N\) is the \(N\)th root of unity.

  This is the same as \(f(x_i) = \sum^{n}_{j=0}a_jx^j_i\) from earlier
  but updated to show that we are using the roots of unity as our
  evaluations points.

  \(DFT(A)\):

  \[a = \begin{bmatrix}1 \\ 2 \\ 0 \\ 0 \end{bmatrix}\]

  \[Va = \begin{bmatrix}
     (1 \cdot 1) + (1 \cdot 2)
     \\ (1 \cdot 1) + (e^{{3\pi\mathrm{i}}/2} \cdot 2)
     \\ (1 \cdot 1) + (e^{3\pi\mathrm{i}} \cdot 2)  
     \\ (1 \cdot 1) + (e^{{9\pi\mathrm{i}}/2} \cdot 2) 
     \end{bmatrix} = 
     \begin{bmatrix} 3 
     \\ 1 + 2e^{{3\pi\mathrm{i}}/2} 
     \\ 1 + 2e^{3\pi\mathrm{i}} 
     \\ 1 + 2e^{{9\pi\mathrm{i}}/2}
     \end{bmatrix}\]

  \(DFT(B)\):

  \[a = \begin{bmatrix}3 \\ 4 \\ 0 \\ 0 \end{bmatrix}\]

  \[Va = \begin{bmatrix}
     (1 \cdot 3) + (1 \cdot 4)
     \\ (1 \cdot 3) + (e^{{3\pi\mathrm{i}}/2} \cdot 4) 
     \\ (1 \cdot 3) + (e^{3\pi\mathrm{i}} \cdot 4) 
     \\ (1 \cdot 3) + (e^{{9\pi\mathrm{i}}/2} \cdot 4) 
     \end{bmatrix} = 
     \begin{bmatrix} 7 
     \\ 3 + 4e^{{3\pi\mathrm{i}}/2}
     \\ 3 + 4e^{3\pi\mathrm{i}} 
     \\ 3 + 4e^{{9\pi\mathrm{i}}/2} 
     \end{bmatrix}\]
\end{itemize}

    \begin{itemize}
\item
  Step 3: Compute the pairwise multiplication \(DFT(A) \cdot DFT(B)\)

  \(DFT(A) \cdot DFT(B)\):
\end{itemize}

\[
     \begin{bmatrix} 3 
     \\ 1 + 2e^{{3\pi\mathrm{i}}/2} 
     \\ 1 + 2e^{3\pi\mathrm{i}} 
     \\ 1 + 2e^{{9\pi\mathrm{i}}/2}
     \end{bmatrix}
     \cdot
     \begin{bmatrix} 7 
     \\ 3 + 4e^{{3\pi\mathrm{i}}/2}
     \\ 3 + 4e^{3\pi\mathrm{i}} 
     \\ 3 + 4e^{{9\pi\mathrm{i}}/2} 
     \end{bmatrix}
     = 
     \begin{bmatrix}
     21
     \\ 3 + 10e^{{3\pi\mathrm{i}}/2} + 8e^{3\pi\mathrm{i}}
     \\ 3 + 10e^{3\pi\mathrm{i}} + 8e^{6\pi\mathrm{i}}
     \\ 3 + 10e^{{9\pi\mathrm{i}}/2} + 8e^{9\pi\mathrm{i}}
     \end{bmatrix}
\]

    \begin{itemize}
\item
  Step 4: Compute the inverse DFT of \(DFT(A) \cdot DFT(B)\).

  This takes us back to coefficient form. The formula for this is
  \[IDFT(DFT(A)\cdot DFT(B)) = x[n] = \frac{1}{N} \sum_{k=0}^{N-1} X[k]\, \omega_N^{-kn}\]
  This simply means we perform the same operation as before but now we
  pick the positive primitive root of unity and pick
  \(a = DFT(A).DFT(B)\), then we divide the resulting \(Va\) by N being
  \(4\) in this case.

  This formula is different from the one we used this process earlier
  because of the percularities of complex numbers and roots of unity but
  it's requires less operations. For instance, we don't need to find the
  inverse of the matrix \(V\).

  \[\omega_4 = e^{{\pi\mathrm{i}}/2}\]

  \[\omega_4 = \{ \omega_4^{0}, \omega_4^{1}, \omega_4^{2}, \omega_4^{3} \} = \{ 1, e^{{\pi\mathrm{i}}/2}, e^{\pi\mathrm{i}}, e^{{3\pi\mathrm{i}}/2} \}\]
\end{itemize}

\[ V = 
     \begin{bmatrix} \omega_4^{0} & \omega_4^{0} & {\omega_4}^{0} & \omega_4^{0}
     \\ \omega_4^{0} & \omega_4^{1} & {\omega_4}^{2} & \omega_4^{3}
     \\ \omega_4^{0} & \omega_4^{2} & {\omega_4}^{4} & \omega_4^{6}
     \\  \omega_4^{0} & \omega_4^{3} & {\omega_4}^{6} & \omega_4^{9}
     \end{bmatrix} = 
     \begin{bmatrix} 1 & 1 & 1 & 1 
     \\ 1 & e^{{\pi\mathrm{i}}/2} & e^{\pi\mathrm{i}} & e^{{3\pi\mathrm{i}}/2}
     \\ 1 & e^{\pi\mathrm{i}} & e^{2\pi\mathrm{i}} & e^{3\pi\mathrm{i}}
     \\ 1 & e^{{3\pi\mathrm{i}}/2} & e^{3\pi\mathrm{i}} & e^{{9\pi\mathrm{i}}/2}
     \end{bmatrix}\]

\[a = \begin{bmatrix}
     21
     \\ 3 + 10e^{{3\pi\mathrm{i}}/2} + 8e^{3\pi\mathrm{i}}
     \\ 3 + 10e^{3\pi\mathrm{i}} + 8e^{6\pi\mathrm{i}}
     \\ 3 + 10e^{{9\pi\mathrm{i}}/2} + 8e^{9\pi\mathrm{i}}
     \end{bmatrix}\]

    \[Va = \begin{bmatrix} 1 & 1 & 1 & 1 
     \\ 1 & e^{{\pi\mathrm{i}}/2} & e^{\pi\mathrm{i}} & e^{{3\pi\mathrm{i}}/2}
     \\ 1 & e^{\pi\mathrm{i}} & e^{2\pi\mathrm{i}} & e^{3\pi\mathrm{i}}
     \\ 1 & e^{{3\pi\mathrm{i}}/2} & e^{3\pi\mathrm{i}} & e^{{9\pi\mathrm{i}}/2}
     \end{bmatrix}
     \cdot
      \begin{bmatrix}
     21
     \\ 3 + 10e^{{3\pi\mathrm{i}}/2} + 8e^{3\pi\mathrm{i}}
     \\ 3 + 10e^{3\pi\mathrm{i}} + 8e^{6\pi\mathrm{i}}
     \\ 3 + 10e^{{9\pi\mathrm{i}}/2} + 8e^{9\pi\mathrm{i}}
     \end{bmatrix}
     =
    \begin{bmatrix}
     21 + (3 + 10e^{{3\pi\mathrm{i}}/2} + 8e^{3\pi\mathrm{i}}) + (3 + 10e^{3\pi\mathrm{i}} + 8e^{6\pi\mathrm{i}}) + (3 + 10e^{{9\pi\mathrm{i}}/2} + 8e^{9\pi\mathrm{i}})
     \\  21 + (e^{{\pi\mathrm{i}}/2} \cdot (3 + 10e^{{3\pi\mathrm{i}}/2} + 8e^{3\pi\mathrm{i}})) + (e^{\pi\mathrm{i}} \cdot (3 + 10e^{3\pi\mathrm{i}} + 8e^{6\pi\mathrm{i}})) + (e^{{3\pi\mathrm{i}}/2} \cdot (3 + 10e^{{9\pi\mathrm{i}}/2} + 8e^{9\pi\mathrm{i}}))
     \\  21 + (e^{\pi\mathrm{i}}  \cdot (3 + 10e^{{3\pi\mathrm{i}}/2} + 8e^{3\pi\mathrm{i}})) + (e^{2\pi\mathrm{i}} \cdot (3 + 10e^{3\pi\mathrm{i}} + 8e^{6\pi\mathrm{i}})) + (e^{3\pi\mathrm{i}} \cdot (3 + 10e^{{9\pi\mathrm{i}}/2} + 8e^{9\pi\mathrm{i}}))
     \\  21 + (e^{{3\pi\mathrm{i}}/2} \cdot (3 + 10e^{{3\pi\mathrm{i}}/2} + 8e^{3\pi\mathrm{i}})) + (e^{3\pi\mathrm{i}} \cdot (3 + 10e^{3\pi\mathrm{i}} + 8e^{6\pi\mathrm{i}})) + (e^{{9\pi\mathrm{i}}/2} \cdot (3 + 10e^{{9\pi\mathrm{i}}/2} + 8e^{9\pi\mathrm{i}}))
     \end{bmatrix}
     =
    \begin{bmatrix}
     12
     \\ 40
     \\ 32
     \\ 0
     \end{bmatrix}
     \]

    \[IDFT(DFT(A)\cdot DFT(B)) = \dfrac{1}{N}Va = \dfrac{1}{4}Va = \dfrac{1}{4} \cdot \begin{bmatrix}
     12
     \\ 40
     \\ 32
     \\ 0
     \end{bmatrix}
     =
     \begin{bmatrix}
     3
     \\ 10
     \\ 8
     \\ 0
     \end{bmatrix}
     \]

    To better understand \(Va\), recall that \(e^{{3\pi\mathrm{i}}/2}\) is
\(-\mathrm{i}\), \(e^{3\pi\mathrm{i}}\) is \(-1\) and
\(e^{{9\pi\mathrm{i}}/2}\) is \(\mathrm{i}\)

As you can see, we have successfully performed polynomial multiplication
using DFT! Finally, let's see how FFT works!

    \subsubsection{How FFT works}\label{how-fft-works}

    FFT is a recursive exploits the structure of the DFT construction and
two properties of roots of unity.

So far, we have explored DFT using matrices and matrix multiplication.
But, to better understand FFT you need to view DFT in terms of the
formulas \(X[k] = \sum_{n=0}^{N-1} x[n]\, \omega_N^{kn}\) and
\(x[n] = \frac{1}{N} \sum_{k=0}^{N-1} X[k]\, \omega_N^{-kn}\) so let's
do that.

Using same \(A(x)\) and \(B(x)\) from above, let's compute \(DFT\) and
\(IDFT\) using the formulas

For DFT: \(X[k] = \sum_{n=0}^{N-1} x[n]\, \omega_N^{kn}\)

\begin{itemize}
\item
  \(A(x)\):

  \begin{itemize}
  \tightlist
  \item
    \(X[0] = x[0] + x[1] + x[2] + x[3] = 1 + 2 + 0 + 0 = 3\)
  \item
    \(X[1] = x[0] + x[1]\omega_4 + x[2]\omega_4^{2} + x[3]\omega_4^{3} = 1 + 2\omega_4\)
  \item
    \(X[2] = x[0] + x[1]\omega_4^{2} + x[2]\omega_4^{4} + x[3]\omega_4^{6} = 1 + 2\omega_4^2\)
  \item
    \(X[3] = x[0] + x[1]\omega_4^{3} + x[2]\omega_4^{6} + x[3]\omega_4^{9} = 1 + 2\omega_4^3\)
  \end{itemize}
\item
  \(B(x)\):

  \begin{itemize}
  \tightlist
  \item
    \(X[0] = x[0] + x[1] + x[2] + x[3] = 3 + 4 + 0 + 0 = 7\)
  \item
    \(X[1] = x[0] + x[1]\omega_4 + x[2]\omega_4^{2} + x[3]\omega_4^{3} = 3 + 4\omega_4\)
  \item
    \(X[2] = x[0] + x[1]\omega_4^{2} + x[2]\omega_4^{4} + x[3]\omega_4^{6} = 3 + 4\omega_4^2\)
  \item
    \(X[3] = x[0] + x[1]\omega_4^{3} + x[2]\omega_4^{6} + x[3]\omega_4^{9} = 3 + 4\omega_4^3\)
  \end{itemize}
\end{itemize}

For IDFT: \(x[n] = \frac{1}{N} \sum_{k=0}^{N-1} X[k]\, \omega_N^{-kn}\)
- \(x[0] = \dfrac{1}{4}(X[0] + X[1] + X[2] + X[3]) = 3\) -
\(x[1] = \dfrac{1}{4}(X[0] + X[1]\omega_4^{-1} + X[2]\omega_4^{-2} + X[3]\omega_4^{-3}) = 10\)
-
\(x[2] = \dfrac{1}{4}(X[0] + X[1]\omega_4^{-2} + X[2]\omega_4^{-4} + X[3]\omega_4^{-6}) = 8\)
-
\(x[3] = \dfrac{1}{4}(X[0] + X[1]\omega_4^{-3} + X[2]\omega_4^{-6} + X[3]\omega_4^{-9}) = 0\)

Secondly, here are the two properties of roots of unity that ultimately
makes FFT possible.

\begin{enumerate}
\def\labelenumi{\arabic{enumi}.}
\tightlist
\item
  \(-\omega_N^k = \omega_N^{k + N/2}\). This is gotten from the fact
  that \(\omega_N^{N/2} = -1\)
\item
  \(\omega_N^2 = \omega_{N/2}\). For example, \(e^{{\pi\mathrm{i}}/4}\)
  is a primitive root from the \(8\)th of unity and
  \((e^{{\pi\mathrm{i}}/4})^2\) is \(e^{{\pi\mathrm{i}}/2}\) being a
  primitive root from the \(4\)th roots of unity.
\end{enumerate}

As you probably guessed, these only holds when \(N\) is even.

    Haven set the foundation, let's look at the FFT algorithm step by step.

Starting from the DFT formula:
\(X[k] = \sum_{n=0}^{N-1} x[n]\, \omega_N^{kn}\)

\begin{itemize}
\item
  \textbf{Step 1}: We are going to split the formula into even/odd
  indices. Setting \(n = 2m\) and \(n = 2m + 1\) so that:
  \[X[k] = \sum_{n=0}^{(N/2)-1} x[2m]\, \omega_N^{k(2m)} + \sum_{n=0}^{(N/2)-1} x[2m + 1]\, \omega_N^{k(2m + 1)} = \sum_{n=0}^{(N/2)-1} x[2m]\, \omega_{N/2}^{km} + \omega_N^{k}\sum_{n=0}^{(N/2)-1} x[2m + 1]\, \omega_{N/2}^{km}\]

  Notice the second property at play here.
\item
  \textbf{Step 2}: Let's label the
  parts:\[E[k] = \sum_{n=0}^{(N/2)-1} x[2m]\, \omega_{N/2}^{km}\]\[O[k] = \sum_{n=0}^{(N/2)-1} x[2m + 1]\, \omega_{N/2}^{km}\]
  So that: \[X[k] = E[k] + \omega_N^{k}O[k]\]
\item
  \textbf{Step 3}: Let's compute: \[X[k] = E[k] + \omega_N^{k}O[k]\]
  \[X[k + N/2] = E[k] + \omega_N^{k + N/2}O[k]\]From the first property,
  we can update the as follows:\[X[k] = E[k] + \omega_N^{k}O[k]\]
  \[X[k + N/2] = E[k] - \omega_N^{k}O[k]\]So, we just have to compute
  \(E[k]\), \(O[k]\) and \(\omega_N^k\) and use it for both
  computations, flipping the sign on the second one.
\item
  \textbf{Step 4}: We perform same steps recursively for \(E[k]\) and
  \(O[k]\) until \(N = 1\).
\end{itemize}


    % Add a bibliography block to the postdoc
    
    
    
\end{document}
